\listfiles
\documentclass[12pt,spanish,letterpaper,oneside]{book}
\usepackage[utf8]{inputenc}
\usepackage[spanish,activeacute,es-lcroman]{babel}	%para incluir los caracteres españoles en el código fuente.
\usepackage{graphicx}
\usepackage{float}		%para que las tablas queden en donde el lugar donde está puesto el código fuente
\usepackage{listings}
\usepackage{color}
\usepackage{url}
\usepackage[breaklinks]{hyperref}
\usepackage{breakurl}	%Con esto por fin se rompieron las líneas del url
\usepackage{listings}	%Para código fuente
\usepackage{fancyvrb}
\usepackage{anysize}	%para establecer los márgenes
\usepackage{fancyhdr}	%Para modificar la posición del número de página
\usepackage{times}
% \renewcommand{\baselinestretch}{1.5} % para saltar 1.5 el espacio

%Título
\title{Dise\~no e implementación de aplicaci\'on cliente en dispositivos iOS para la reproducci\'on de medios en un servidor de streaming timeshift.}
          
%Información del autor
\author{
	Javier Cristóbal González Ovalle									\\
	Ingeniería Civil Electrónica								\\
	Universidad Técnica Federico Santa María		\\	
	email:\url{javier.gonzalezo@alumnos.utfsm.cl}	\\
																		\\
	Profesor Guía: Agustín González Valenzuela	\\
	Profesor Co-referente: Tomás Arredondo Vidal		\\
} 
%Márgenes
\marginsize{4cm}{3cm}{3cm}{3cm}
%Fecha	
\date{Agosto de 2012}
\sloppy


%Colocar el número de página abajo a la derecha

%Configuración de las páginas planas.(Con esto logré la posición del número
%en la tabla de contenidos, primera página de capítulo, etc.
 \pagestyle{fancyplain}
 \fancyhf{}
 \rfoot{\thepage}
 
 %\pagenumbering{roman} 
%Estructura del documento
\begin{document}
	%\frontmatter
	%Incluir Portada
	\begin{titlepage}

%Cabecera de presentación de la universidad.
\begin{center}

	\large  UNIVERSIDAD TÉCNICA FEDERICO SANTA MARÍA 	\\
 	\large \sc DEPARTAMENTO DE ELECTRÓNICA						\\
 	\large \sc VALPARAÍSO-CHILE												\\
\begin{figure}[h!]
	\centering
 	\includegraphics[scale=1]{imgs/logo_utfsm.eps} 
\end{figure}
\end{center}

%Título de la memoria
\begin{center}
 \LARGE{Diseño e implementación de aplicación cliente en dispositivos iOS para la reproducción de medios en un servidor de streaming timeshift.}
\end{center}

%Información del alumnos
\begin{center}
 \normalsize Javier Cristóbal González Ovalle\\
 \normalsize MEMORIA DE TITULACIÓN PARA OPTAR AL TÍTULO DE
 \normalsize INGENIERO CIVIL ELECTRÓNICO			\\
 \normalsize MENCIÓN COMPUTADORES
\end{center}
\vspace*{1.5cm}

\begin{center}
\begin{tabular}{ll}
 \normalsize Profesor Guía: & Agustín González Valenzuela\\
 \normalsize Profesores Co-Referentes: &Tomás Arredondo Vidal\\
 															&Alfredo Luco Díaz
\end{tabular}
\end{center}
 
\vspace*{2cm}
\begin{center}
 \normalsize Valparaíso, Agosto de 2012. \\
 %\normalsize Material de referencia, su uso no involucra responsabilidad del autor o de la Institución.
\end{center}
\end{titlepage}
	%\begin{titlepage}

%Cabecera de presentación de la universidad.
\begin{center}

	\large  UNIVERSIDAD TÉCNICA FEDERICO SANTA MARÍA 	\\
 	\large \sc DEPARTAMENTO DE ELECTRÓNICA						\\
 	\large \sc VALPARAÍSO-CHILE												\\
\begin{figure}[h!]
	\centering
 	\includegraphics[scale=1]{imgs/logo_utfsm.eps} 
\end{figure}
\end{center}

%Título de la memoria
\begin{center}
 \LARGE{Diseño e implementación de aplicación cliente en dispositivos iOS para la reproducción de medios en un servidor de streaming timeshift.}
\end{center}

%Información del alumnos
\begin{center}
 \normalsize Javier Cristóbal González Ovalle\\
 \normalsize MEMORIA DE TITULACIÓN PARA OPTAR AL TÍTULO DE
 \normalsize INGENIERO CIVIL ELECTRÓNICO			\\
 \normalsize MENCIÓN COMPUTADORES
\end{center}
\vspace*{1.5cm}

\begin{center}
\begin{tabular}{ll}
 \normalsize Profesor Guía: & Agustín González Valenzuela\\
 \normalsize Profesores Co-Referentes: &Tomás Arredondo Vidal\\
 															&Alfredo Luco Díaz
\end{tabular}
\end{center}
 
\vspace*{2cm}
\begin{center}
 \normalsize Valparaíso, Agosto de 2012. \\
 %\normalsize Material de referencia, su uso no involucra responsabilidad del autor o de la Institución.
\end{center}
\end{titlepage} si le pongo include no me aparece con la compilación rápida
	
	%Comenzar numeración romana
	\renewcommand{\thepage}{\roman{page}}
 	\setcounter{page}{1}
	
	%Índice General
	\tableofcontents
	\addcontentsline{toc}{chapter}{Índice General}

	%Índice de Figuras 
 	\listoffigures
 	\addcontentsline{toc}{chapter}{Índice de Figuras}

	%Índice de Tablas
 	\listoftables
 	\addcontentsline{toc}{chapter}{Índice de Tablas} 
 
   \newpage
\thispagestyle{empty}
\begin{center}
 \Large \textbf{Diseño e implementación de aplicación cliente en dispositivos iOS para la reproducción de medios en un servidor de streaming timeshift}\\

\normalsize Memoria para optar al título de Ingeniero Civil Electrónico, mención Computadores \\
\normalsize Javier González Ovalle \\
\normalsize Profesor Guía: Agustín González Valenzuela \\
\normalsize Junio de 2012

\Large \textbf{Resumen}

\end{center}
\normalsize
%\textbf{Objetivos del Trabajo:} \\
\subsection*{Objetivos del Trabajo}
\normalsize
%Siguiendo el esquema de AHidalgo.pdf

%Primer parrafo, descripción breve de los objetivos
Nuevos avances se han realizado en el manejo de transmisiones en vivo de medios como audio y/o video, que permiten emular la característica principal de dispositivos de sobremesa del tipo \textquotedblleft Programmable (Digital) Video Recorders\textquotedblright \ (PVR ó DVR), los cuales permite pausar, retroceder o avanzar, cuando es posible, la linea de tiempo de reproducción. Esto se logra gracias a aplicaciones que interactuan con servidores de flujo de datos de audio y video para lograr un registro de la reproducción, entregando al cliente comunicación eficaz para recibir datos y retornar comandos.\\ 

%Segundo párrafo, contextualizaciòn del problema
Se busca suplir la necesidad de un reproductor de audio y/o video que permita la interacción con un servidor de \textit{streaming} (flujo) de datos, con habilidad de \textit{Time-Shift} (cambio temporal). Esto es la habilidad de retroceder o avanzar (si es posible) en la linea de tiempo de la reproducción del audio o del video al que el usuario atiende.\\

Esto ya se logra en computadores de escritorio gracias a un reproductor de streaming basado en la tecnolog\'ia Flash de Adobe. Los reproductores embebidos en p\'aginas web reproducen el stream en base al protocolo Real Time Messaging Protocol (RTMP), el cual permite llamadas de control entre cliente y servidor para manejar la linea de tiempo. Sin embargo se busca llegar a un mayor público portando la capacidad de reproducci\'on en dispositivos que utilizan el sistema operativo iOS de Apple.\\

El objetivo de este trabajo se debe a que Apple por políticas internas, no apoya la distribuci\'on de Adobe Flash. Por lo tanto se busca implementar una soluci\'on que permita la reproducci\'on de la transmisi\'on con \textit{Time-Shift} y que adem\'as concuerde con las pol\'iticas de distribuci\'on para las aplicaciones iOS en la AppStore.
 
\large
%\textbf{Trabajo a Desarrollar:}\\
\subsection*{Trabajo a Desarrollar}
\normalsize
%Tercer y cuarto párrafo, que y como se realizó
El trabajo toma como punto de partida la investigaci\'on de los protocolos propuestos por Apple para la distribuci\'on de medios (audio y/o video). En base a estos protocolos buscar similitudes con el sistema de streaming para computadores de escritorio. Adem\'as se busca implementar un m\'etodo de marcaci\'on de tiempo en el stream de datos destinado a dispositivos m\'oviles, de forma que el servidor los env\'ie para que el cliente los interprete. Este cliente debe consistir en una aplicación que reciba y env\'ie comandos compatibles con el servidor Wowza para controlar la informaci\'on del medio a reproducir. Finalmente se busca diseñar una interfaz gr\'afica que represente la disponibilidad de tiempo y permita controlar la reproducción.\\

%\textbf{Evaluaciones a Realizar:}
\subsubsection*{Evaluaciones a Realizar}
\normalsize
\begin{itemize}
\item	Comprobar la transmisión y recepción de datos en HTTP Live Stream u otro protocolo compatible con iOS, a través de un servidor de streaming dedicado y un dispositivo con iOS.
\item	Comprobar el intercambio de mensajes de control a través de un enlace entre el cliente móvil y el servidor de streaming.
\item	Pruebas comparativas de la reproducción en dispositivo móviles y computadores de escritorio.
\end{itemize}

%\textbf{Resultados esperados:}\\
\subsubsection*{Resultados esperados}
\normalsize
Se pretende que la reproducción en dispositivos móviles de las transmisiones en vivo se asemejen a la experiencia de reproducción en computadores de escritorio.\\

%Palabras clave.
%	\textbf{Palabras Claves:} streaming, iOS, video, audio, apple, iPhone.

   \addcontentsline{toc}{chapter}{Resumen}

   % \newpage
%\thispagestyle{empty}
\begin{center}
% Diseño e implementación de aplicación cliente en dispositivos iOS para la reproducción de medios en un servidor de streaming timeshift
 \Large \textbf{Design and Implementation of a Client iOS Application for a Streaming Media Server with Timeshift Capabilities}\\
  

%\normalsize Memoria (buscar traducción de memoria) para optar al título de Ingeniero Civil Electrónico, mención Computadores \\
\normalsize Javier Cristóbal González Ovalle \\
%\normalsize Profesor : Agustín González Valenzuela \\
\normalsize August 2013

\Large \textbf{Abstract}

\end{center}
\normalsize

%Primer parrafo, descripción breve de los objetivos
The goal of this project is to extend for mobile devices the streaming media services of the company AltaVoz S.A. The main feature of their services is to give the audience a certain control over what they are watching, by being able to seek to another time of playback. This capability gives to the user the freedom of watching past shows, repeat highlights, or pause the moment for later playback. The feature is being promoted as \textit{timeshift}.\\



%Segundo párrafo, contextualizaciòn del problema
The \textit{timeshift} streaming media service is currently available for the Web platform, exclusively for computers able to display Adobe Flash content. The main issue for this service comes from the interest of potential clients which are focused on broadcasting their content to mobile devices, specifically smartphones capable of media playback. Since a great amount of devices are not able to load the Adobe Flash plug-in in their Web Browsers, there is a necessity for an alternative method of media playback with the \textit{timeshift} feature.\\

 
%Tercer y cuarto párrafo, que y como se realizó


Research started with the current protocols being used at the moment for the service developed by AltaVoz S.A., these were compared with the options for streaming media playback given by Apple Inc. for their iOS mobile operating system, the results meant several points which were needed to be modified in the server side, this way the server would deliver compatible streaming media to iOS devices. Secondly was the development of an iOS application able to play the content delivered by the modified \textit{timeshift} server, also featuring a simple user interface for the \textit{timeshift} capability.
Later was needed the implementation of a sharing system for time stamps, this was succeed by integrating the social network Twitter. To conclude it was necessary to check the behavior for different scenarios of Internet connectivity in the mobile device. 
\\

 
   % \addcontentsline{toc}{chapter}{Abstract} 
 
 	\mainmatter
	%Reiniciar numeración y cambiar a números áraes
	\renewcommand{\thepage}{\arabic{page}} 
	\setcounter{page}{1}
	
	%Configuración del estilo de las páginas generales.(Como el cuerpo)
	\pagestyle{fancy}
	%Colocar el número de página abajo a la derecha
    \fancyhf{}
    \rfoot{\thepage}
    \lhead{\leftmark} 	
	\thispagestyle{fancy}
 	
	%Incluir capítulo 1
	%capitulo1.tex
\chapter{Introducción}
%Introducción nueva
%¿Qué son las redes de sensores inalámbricos?
Tomando en cuenta que desde sus inicios la reproducción de medios siempre ha estado relacionada con los últimos avances científicos, los medios de distribución actuales tienen sus bases en los avances que datan  del siglo XIX.\\

Una de las principales aplicaciones de estas redes es la construcción de sistemas de monitoreo que recolectan mediciones de variables físicas en zonas donde el monitoreo cableado representa costos muy elevados o no es posible por la falta de redes de distribución eléctrica\cite{WileyWSN}.\\
%¿Cuáles son las áreas de investigación y desarrollo de WSN?
%Las principales temas de investigación (medidas en términos de frecuencia de publicaciones) son: Despliegue, seguimiento de objetivos, localización, recolección de datos, ruteo y agregación, seguridad, protocolos de acceso al medio (MAC), acceso a bases de datos y sincronización del tiempo. Estos temas son tratados en más del cincuenta porciento del total de publicaciones\cite{WileyWSN} .\\ %buscar como continuar después de una itemezación con dos puntos y separados por comas.

%Presentar el tema del ruteo y recolección de datos
Dentro de un sistema de recolección de datos, basado en una red de sensores inalámbricos, es necesario contar con un protocolo de ruteo que permita direccionar los datos hacia uno o varios puntos de recolección dentro de la red. Dadas las características y limitaciones de los sensores inalámbricos, el principal objetivo de estos protocolos es reducir el consumo de energía de la red, y a la vez, mantener su funcionalidad frente a las diferentes fallas que se puedan presentar\cite{WSNSurvey}.\\

%IR DE LO GENERAL A LO PARTICULAR!
%El capítulo de introducción debe comenzar explicando que son las 	WSNs.
\section{Redes de sensores inalámbricos}
%Cosas que hay que dejar claras:
%	-Es un sistema distribuído.	OK
%	-Como nacen las WNSs.
%	-Formada por motes.	OK
%	-Qué cosas permiten las WSNs y por qué las permiten
Una red de sensores inalámbricos es un sistema distribuido formado por sensores autónomos desplegados espacialmente, que en forma cooperativa, reúnen información del ambiente que los rodea. Los dispositivos utilizados para formar estas redes son conocidos como ''motes'' y poseen capacidad de procesamiento, capacidades sensoriales, capacidad de comunicación inalámbrica y son autónomos energéticamente.

\section{Aplicaciones}
%En el mundo
De los trabajos publicados se está en conocimiento de que las redes de sensores se están aplicando en\cite{WileyWSN}: \\

\textbf{Alerta de desastres}: La aplicación más recurrente en esta área es la detección de incendios, donde los dispositivos utilizan sensores de temperatura y algún sistema de localización para alertar en caso de incendio y/o entregar a los servicios de emergencia información relacionada con la distribución del incendio.\\

\textbf{Edificios inteligentes}: Principalmente, el uso de redes de sensores inalámbricos para entregar información a los sistemas de control de aire acondicionado y ventilación. También permiten la reducción de costos frente a los sistemas actuales que funcionan en su mayoría utilizando dispositivos cableados. También se han propuesto sistemas para monitorear los niveles de estrés mecánico en la estructura de los edificios construidos en territorios de alta actividad sísmica.\\

\textbf{Agricultura de precisión}: Las redes de sensores inalámbricas aplicadas a la agricultura están, generalmente, construidas para realizar mediciones de suelo: humedad, temperatura y sales. Se ha determinado que un sensor por hectárea es suficiente para obtener datos necesarios para monitorear el estado de los cultivos.\\

%En chile

%Red para monitoreo forestal
En Chile, durante los años 2006 y 2007 se desarrolló el proyecto ''Redes inalámbricas de sensores para la detección de incendios forestales y monitoreo de variables de estado de combustibles para el recurso forestal de la Región de Valparaíso''\cite{ProyectoForestal}. El proyecto fue desarrollado en conjunto por ADDERE Ltda., CONAF, Universidad de Concepción, Forestal ACE Internacional S.A., ONEMI, CORMA y la Agencia Chilena del Espacio. El objetivo de este proyecto era monitorear el estado de los combustibles forestales y detectar los posibles focos de incendio en tiempo real, debido a que anualmente se queman más de seis mil hectáreas sólo en la Región de Valparaíso, causando pérdidas materiales por más de cincuenta millones de dólares.\\

%Wiseconn
Wiseconn es una empresa chilena que se ha dedicado a desarrollar productos para la agricultura de precisión utilizando redes de sensores inalámbricos. Actualmente utilizan tecnología basada en la especificación Zigbee para recolectar datos relacionados con mediciones de suelo \cite{Wiseconn}.

%Austec

\section{Protocolo de ruteo para un sistema de recolección de datos}
%\section{Visión general de los requerimientos}
%Hablar de la aplicación objetivo y cuáles son sus requerimientos 
%La implementación de un protocolo de ruteo se realiza como parte de un sistema de recolección de datos orientado a. Por ejemplo, el sistema se puede desplegar en un terreno de cultivo, donde se desee principalmente medir la temperatura y la humedad del suelo, luego concentrar esos datos en un Gateway que enviaría esos datos a través de una red ethernet a un sistema computacional que almacena, analiza y presenta esos datos a través de una aplicación web. \\

%Títulos alternativos
El objetivo de este trabajo es implementar un protocolo de ruteo multihop como parte de un sistema de recolección de datos que está construido sobre una red de sensores inalámbricos. Este sistema está orientado principalmente al monitoreo de variables físicas principalmente relacionadas con la agricultura, como la humedad y la temperatura de los suelos, en zonas de difícil acceso y que no cuentan con redes de distribución eléctrica.

\begin{figure}[H]
% \includegraphics[scale=0.625]{imgs/ArquitecturaGeneral.eps} 
\caption{Arquitectura general del sistema de recolección.}
\end{figure}

En la figura 1.1, se describe la arquitectura del sistema de recolección de datos. Este sistema tiene como objetivo principal: mover los datos desde los puntos en la red donde se realizan las mediciones, hasta el punto en la red donde los datos son concentrados y enviados por otro medio de transmisión en el camino a su destino final. Para realizar esta tarea se debe implementar una capa de red, que se encargue de controlar la topología de la red y establecer rutas para los mensajes que transportan los datos. El sistema está ideado para utilizar una red de entre 30 y 50 nodos, donde todos los nodos (excluyendo el Gateway) son energéticamente autónomos y su tiempo de operación debe alcanzar el orden de meses.

Los requerimientos establecidos para la implementación del protocolo son los siguientes:

\begin{enumerate}
\item \textbf{Rutear los mensajes hacia el Gateway:} Se debe implementar un protocolo que realice el ''mejor esfuerzo'' en transportar los mensajes desde un nodo origen, reenviar los mensajes por nodos intermedios y entregarlos al Gateway.Se debe contemplar que el nodo puede recibir mensajes de múltiples nodos en un periodo de tiempo reducido, por lo que debe de alguna manera almacenar los mensajes recibidos y enviarlos en orden, minimizando la pérdida de mensajes.\\

\item \textbf{Tolerar falla de desaparición de nodos:} En el contexto de la comunicación multi-hop son muchas las fallas e inconsistencias que se pueden generar, sin embargo, el protocolo multi-hop debe al menos tolerar la desaparición de nodos de la red, es decir, si un nodo es sacado de la red o agota su batería, entonces la red no debe dejar de funcionar y los nodos que se queden sin ruta, deben buscar otra si es que existe.\\

\item\textbf{Reducir el consumo energético:} Dado que la principal variable de diseño, en todos los aspectos, de una red de sensores inalámbricos es el consumo energético, el protocolo implementado debe realizar operaciones que permitan gestionar el consumo y alcanzar un tiempo de operación del orden de meses. \\
\end{enumerate}

En el marco de este trabajo, los objetivos generales son los siguientes:
\begin{enumerate}
  \item Informar sobre los protocolos de ruteo de mensajes para redes de sensores inalámbricos.
  \item Desarrollar una red de sensores inalámbricos, utilizando la plataforma de hardware TmoteSky.
  \item Desarrollar una aplicación para la recolección de datos de la red.
  \item Evaluar el consumo energético de la red implementada.
\end{enumerate}

 El primer objetivo se desarrolla en profundidad en el capítulo 2. El segundo objetivo se desarrolla en los capítulos 3, 4 y 5, donde se plantea un diseño de protocolo y posteriormente como se implementa éste. El tercer objetivo se desarrolla en los capítulos 3 y 5 como parte de las verificaciones del sistema. Finalmente, en el capítulo 5 se evalúa el consumo energético el protocolo.\\
 
\section{Estado del Arte}
El estado del arte para un sistema de recolección de datos se describe utilizando la siguiente estructura: primero, se revisan los algoritmos más recurrentes en la literatura orientados a la recolección de datos; segundo, se describen los algoritmos implementados para las plataformas más populares; finalmente, se describe el software y hardware de desarrollo disponible.

\subsection{Protocolo de ruteo}
Según el modelo OSI para el diseño de comunicación entre computadores, la capa de red es responsable de todo lo relacionado con entrega de paquetes de un computador origen a uno destino. Esto incluye el direccionamiento lógico de mensajes y el ruteo de paquetes a través de una o varias redes. Un componente de la implementación de la capa de red son los algoritmos de ruteo. \\

%\begin{figure}[H]
 %\centering
 %\includegraphics[scale=0.8]{imgs/PilaModeloOSI.eps}
 %\caption{Pila de red establecida por el Modelo OSI}
%\end{figure}

Los algoritmos de ruteo, son métodos que permiten definir rutas dentro de una red para transmitir paquetes de datos. La estrategia de ruteo elegida puede tener un impacto significativo en el desempeño de la red\cite{WSNSurvey}.\\

%Las redes de sensores inalámbricos son extremadamente versátiles, por lo que la forma en que éstas son desplegadas, dependerá de la naturaleza de la aplicación. Las redes pueden estar compuestas de nodos estáticos o móviles. Por ejemplo: para aplicaciones de monitoreo, los sensores son desplegados en un área y, usualmente, se mantienen estáticos mientras recolectan información; para aplicaciones en el área de la salud, los nodos pueden ser adheridos a insumos o personas para monitorear sus signos vitales, en este caso, la red debe permitir la movilidad de los nodos.\\

Un primer enfoque para diseminar datos en la red, es enviar los datos recolectados por un nodo en un salto (comunicación directa) hacia el gateway, sin embargo, si el nodo está muy lejos del gateway, puede que éste nunca reciba el mensaje. Por lo tanto, el siguiente enfoque son los protocolos multi-hop, donde varios nodos intermedios se encargan de llevar un paquete de datos desde fuente a destino. En un algoritmo multi-hop, cada nodo ejecuta el algoritmo de ruteo que determinará cuál es el próximo destino de un paquete en la red.

\subsection{Algoritmos de ruteo propuestos en la literatura}
Se han propuesto diversos algoritmos de ruteo en la literatura, sin embargo, en este trabajo se describirán en detalle los algoritmos de ruteo jerárquicos, que según \cite{WSNSurvey} son los más indicados para aplicaciones de recolección de bajo consumo. Dentro de los algoritmos de ruteo jerárquico, los más recurrentes son: LEACH, PEGASIS y TEEN. Se profundizará en estos protocolos en el siguiente capítulo.

%En esta sección debería dar una breve reseña  de cada protocolo, con sus referencias.

\subsection{Protocolos implementados}
En los últimos años, se han implementado diversas capas de red para distintos sistemas.  Cada implementación de capa de red, está asociada a un protocolo de red en particular. En el sistema TinyOS, existen dos implementaciones de red que son distribuidas junto con el sistema: Dissemination (Diseminación) y Collection Tree Protocol (Protocolo de árbol de recolección). Por otro lado, Texas Instruments ha implementado dos tipos de red para su hardware: Z-Stack, la implementación de Zigbee de TI; y SimpliciTI, un protocolo de red simple para construir redes de topología estrella. En el siguiente capítulo, se realizará una descripción detallada de estos protocolos.

\subsection{Hardware de desarrollo}
%De lo general a lo particular
% Qué tipo de hardware existe?
En la actualidad, existe una gran variedad de hardware que puede ser utilizado para el desarrollo e implementación de una red de sensores inalámbricos. Como se ha mencionado antes, un sensor inalámbrico (mote) está compuesto al menos por: Un microcontrolador, que permite administrar cada recurso del nodo; un transceptor, que permite comunicar el nodo con otros dispositivos; sensores y actuadores, que permiten al nodo interactuar con el entorno; batería, si es que el nodo está diseñado para que opere de manera autónoma. Los principales fabricantes del hardware para estas redes son: Texas Instruments, Nordic Semiconductors, Atmel y Jennic.\\

%Debo incluir info sobre los motes que soportan tinyos
\subsubsection{Plataformas TinyOS}
Algunas de las plataformas de desarrollo soportadas actualmente por TinyOS 2.1 son:

\begin{center}
\begin{table}[H]
\caption{Algunas plataformas de hardware soportadas por TinyOS} 
\begin{tabular}{|l|p{10cm}|}
\hline \textbf{Plataforma} & \textbf{Descripción} \\ 

\hline Intelmote2 \cite{IntelMote2Web} &  Plataforma basada en el procesador PXA271 XScale, tranceptor CC2420, posee dos conectores de 31 y 21 pines para sensores,  e interfaz USB. Utiliza dos baterías AA\\ 

%\hline Iris \cite{IrisDatasheet} & Plataforma basada en el SoC XM2110CA de Atmel, conector de expansión de 51 pines para sensores. Utiliza dos baterías AA \\ 

\hline Mica2 \cite{Mica2Datasheet} &  Plataforma basada en el Soc de bajo consumo Atmel ATmega 128L, conector de expansión de 51 pines para sensores. Utiliza dos baterías AA. \\ 

\hline Micaz \cite{MicazDatasheet} &  Plataforma basada en el Soc de bajo consumo Atmel ATmega 128L, conector de expansión de 51 pines para sensores. \\ 

%\hline Shimmer \cite{ShimmerWiki} &  Plataforma diseñada para vestir, basada en el microprocesador de bajo consumo MSP430F1611. Incorpora un tranceptor Bluetooth WML-C46A y un tranceptor CC2420. Posee un lector de tarjetas Micro-SD para almacenamiento masivo. Incluye un acelerómetro de 3-ejes. Tiene una batería recargable de Ion-Litio integrada.\\ 

\hline TelosB (TmoteSky) \cite{TelosBDatasheet} & Plataforma basada en el microcontrolador de bajo consumo MSP430, utiliza el transceptor CC2420, incluye sensores de radiación solar total y radiación fotosintétizable. Utiliza dos pilas AA.  \\ 

%\hline Tinynode \cite{TinyNodeWeb} &  Plataforma basada en el microcontrolador de bajo consumo MSP430, utiliza el tranceptor de bajo consumo Xemics  XE 1205, conector de expansión de 31 pines para conectar sensores e interfaces de comunicación serial.\\ 
\hline 

\end{tabular}

\end{table} 
\end{center}
%Se omiten Mica, Mica2dot, eyesIFX  y telos a porque parece que están descontinuados


\subsubsection{Plataformas Zigbee}
Texas Instruments es uno de los principales impulsores de esta tecnología, ellos comercializan las siguientes plataformas de desarrollo.
%Debo incluir info sobre los motes que soportan Zigbee, simpliciti.
\begin{table}[H]
\centering
\caption{Plataformas de desarrollo Zigbee}
\begin{tabular}{|l|p{10cm}|}
\hline \textbf{Plataforma} & \textbf{Descripción} \\ 
\hline CC2430ZDK & Kit de desarrollo para redes Zigbee que incluye tarjetas de desarrollo para el SoC CC2430. Poseen interfaz USB y utiliza dos baterías AA. \\ 
\hline CC2530ZDK & Kit de desarrollo para redes Zigbee que incluye tarjetas de desarrollo para el SoC CC2530. Poseen interfaz USB y utiliza dos baterías AA.\\ 
\hline 
\end{tabular} 
\end{table}

\subsubsection{Otras plataformas}
A continuación se describen otras plataformas para redes de sensores inalámbricos.
\begin{table}[H]
\centering
\caption{Otras plataformas de desarrollo para redes de sensores inalámbricos}
\begin{tabular}{|l|p{10cm}|}
\hline \textbf{Plataforma} & \textbf{Descripción} \\ 
%\hline ATAVRRZ201 \cite{ATAVRRZ201}& Kit de desarrollo para redes Zigbee basado en el microcontrolador ATMega1281V y el tranceptor AT86RF230.\\ 
\hline AVR Raven & Kit de demostración para redes Zigbee que incluye un dispositivo que integra tres chips: el ATmega3290P para manejar los dispositivos I/O; el ATMega1284P para ejecutar la pila de comunicación; y la radio AT86RF230. Estas tarjetas incluyen un Display, un parlante, micrófono, memoria flash, reloj de tiempo real, botones y switchs. Utiliza una batería de tipo botón o alimentación externa desde 5V hasta 12V.\\
\hline EM430F6137RF900 & Plataforma basada en el SoC CC430, que integra un MSP430 y una radio CC1100 en un chip. Incluye un conector para Display LCD que se distribuye junto con la plataforma. Puede utilizar dos baterías AAA. \\ 
\hline ez430 Chronos & Plataforma de desarrollo embebida en un reloj deportivo. Esta plataforma está basada en el SoC CC2430, e incluye acelerómetro y sensor de presión, además de los de temperatura y voltaje del CC430. Utiliza una batería de celda tipo botón.\\ 
\hline SunSpot & Los dispositivos SunSPOTs son motes para redes de sensores inalámbricos desarrollados por Sun Microsystems. Estos dispositvos utilizan el estándar IEEE 802.15.4 para comunicación inalámbrica y tienen integrada la máquina virtual de Java (Squawk). Además, cada SunSPOT cuenta con acelerómetro, sensor de temperatura, sensor de luz, 8 leds tri-color y 2 switches. \\ 
\hline 
\end{tabular} 
\end{table}


\subsection{Software para redes de sensores inalámbricos}
%TinyOS, ZStack, SimpliciTI, etc
%Breve descripción de tinyOS e indicación para revisar el anexo XX.
\subsubsection{TinyOS}
TinyOS es un sistema operativo embebido de código abierto, diseñado especialmente para desarrollar aplicaciones para redes de sensores inalámbricos. Su desarrollo comenzó en la Universidad de Berkeley a fines de los años 90 y actualmente su desarrollo está a cargo de la TinyOS Alliance. TinyOS está diseñado para permitir concurrencia, operación con recursos limitados, soporte de un rango amplio de aplicaciones y hardware, y adaptación a la evolución del hardware. TinyOS es un sistema conducido por eventos  programado en NesC, un dialecto derivado de C que reduce el consumo de memoria, utiliza la programación orientada a componentes y provee la detección de ``race conditions''. \\

\subsubsection{Contiki}
Contiki es otro sistema operativo embebido de código abierto, diseñado para ser portable, multi-tarea, y con un consumo de memoria reducido. Contiki es desarrollado en el Swedish Instute of Computer Science y actualmente cuenta con desarrolladores de la industria como Cisco y Atmel. Este sistema provee comunicación vía IPv4 e IPv6, mediante la utilización de una versión reducida de la pila IP diseñada para microcontroladores. Además, provee un sistema de archivos para almacenar información en la memoria flash de los nodos. Actualmente, Contiki soporta el hardware TmoteSky (Telosb).

%Breve descripción de ZStack, basarse en pagina de ti del zstack
\subsubsection{Z-Stack}
Z-Stack es la implementación del standard Zigbee desarrollada por Texas Instruments para su portafolio de hardware compatible con el IEEE 802.15.4. La versión actual implementa las especificaciones Zigbee 2007, Zigbee Pro y Zigbee 2006. Los chips que soportan este software son: CC2530, MSP430+CC2520, CC2430 y MSP430+CC2420. 	Esta implementación se construye sobre el sistema operativo OSAL que provee herramientas como Tasks (Tareas), Timers y un manejo más amigable de las interrupciones. Es importante destacar que ninguna de las versiones implementa la topología Cluster-Tree para una red completa de sensores independientes energéticamente que propone la especificación.

%Breve descripción de SimpliciTI, basarse en pagina de ti de simpliciti
\subsubsection{SimpliciTI}
SimpliciTI es un protocolo de bajo consumo desarrollado por Texas Instruments 	para sus productos Sub-1GHz y portafolio compatible con el IEEE 802.15.4. Es un protocolo reducido que permite implementar redes con topología estrella, donde dispositivos always-on  coordinan la red. A diferencia de Z-Stack, SimpliciTI no está construido sobre un sistema operativo, por lo que no provee abstracciones de hardware ni herramientas de programación como OSAL.	

%Breve descripción de la de los sunspots, basarse en documentacion
\subsubsection{Java y SunSpot}
Los SunSPOTs son distribuidos junto a un kit de desarrollo que trae las herramientas, tutoriales y ejemplos necesarios para realizar la programación de estos dispositivos. El kit incluye el IDE Netbeans y todos los accesorios (compiladores, depuradores, interfaz de software para comunicación USB) incorporados al IDE. Estos dispositivos son programados en JAVA, utilizando la máquina virtual Squawk para ejecutar el bytecode en el dispositivo. La máquina virtual java Squawk es un desarrollo experimental de código abierto con el objetivo de acelerar la portabilidad y funcionalidad de la máquina virtual Java. Las principales características que permiten su implementación para un dispositivo embebido, son el tamaño reducido del bytecode generado (35\%  a 45\% de una clase equivalente en J2ME)  y una arquitectura que permite ahorrar memoria del dispositivo.

%Breve descripción de alguna otra de otro fabricante..freescale o atmel.






	
	%Incluir capítulo 2
	\chapter{Investigación sobre protocolos de ruteo}
%Breve introducción
La primera etapa del proyecto consiste en investigar los algoritmos de ruteo que se han propuesto en la literatura, explorando la clasificación de estos y describiendo en detalle los principales algoritmos jerárquicos. Posteriormente, se describen los protocolos que provee el sistema TinyOS (el que utiliza el TMoteSky, y sobre el cuál se desarrollarán las siguientes etapas de este trabajo) y los que especifica el estándar de red Zigbee. Al final el capítulo, se revisan los resultados obtenidos por algoritmos que se han propuesto en la literatura y la selección de uno de ellos para su implementación.

\section{Algoritmos de ruteo propuestos en la literatura}
%Introducción
En la literatura se han propuesto diversos algoritmos de ruteo y que se clasifican de distintas maneras. La clasificación más recurrente se realiza '' por topología de la red '' \cite{WSNSurvey} que agrupa los protocolos como: ''flat routing'', ''Hierarchical routing'' y ''Geo routing''.\\

\textbf{Flat routing:} En este tipo de algoritmos, todos los nodos tienen el mismo rol en el proceso de ruteo participando en el reenvío de los mensajes. Se asume que el tráfico ocurre de manera aleatoria y direccionado a cualquier punto de la red. Estos protocolos tienen la ventaja de balancear la carga en la red y ser flexibles en cuanto a tolerar las fallas o cambios en la red.\\

\textbf{Hierachical routing:} Estos algoritmos utilizan estructuras jerárquicas, donde los nodos asumen distintos roles en el proceso de ruteo. Estos protocolos están orientados a la recolección, donde el tráfico se genera de manera aleatoria o sincronizada, y direccionado a uno o varios puntos de recolección.\\

\textbf{Geographical Routing:} Estos algoritmos utilizan la posición geográfica del nodo para realizar el proceso de ruteo. Usualmente estos protocolos se dividen en dos sub categorías: ''Basados en posición'', donde se realiza el ruteo de mensajes unicast (dirigidos a un nodo en particular) en función de la posición geográfica o relativa de un nodo; y los del tipo ''Geocasting'' donde los mensajes son ruteados hacia grandes grupos de nodos, en función de la posición del grupo.\\

Para sistemas de recolección de datos que necesitan operaciones de bajo consumo, se indica que los algoritmos jerárquicos son los más adecuados, ya que están diseñados para este tipo de escenarios \cite{WSNSurvey}. A continuación, en las secciones 2.1.1, 2.1.2 y 2.1.3, se describirán los algoritmos jerárquicos más recurrentes en la literatura.

%Descripción de algunos algoritmos
%Protocolo LEACH
\subsection{Low Energy Adaptive Clustering Hierarchy}
%¿Qué es LEACH?
Low Energy Adaptive Clustering Hierarchy (LEACH) \cite{LeachOriginal}es un algoritmo para redes de sensores inalámbricos, diseñado para: extender el tiempo de autonomía de la red, mediante la distribución de las tareas de red de manera homogénea; y reducir el tráfico de mensajes hacia la estación base, mediante el uso de agregación de datos.\\

%En general, ¿Cómo funciona?
LEACH organiza la red en clusters, donde cada cluster está formado por un ``Cluster Head'' y ``Cluster Members''. En cada cluster, el Cluster Head es el encargado de recibir y reenviar los mensajes generados por los Clusters Member. Estos roles están asociados a distintos consumos energéticos según el modelo energético que se describe más adelante, debido a la distinta cantidad de mensajes que deben manejar. LEACH utiliza la rotación aleatoria del rol de Cluster Head para distribuir de manera homogénea el consumo energético dentro de la red de sensores. Para permitir escalabilidad y robustez, utiliza coordinación local de los nodos, permitiendo una red dinámica. Además, sugiere realizar agregación de datos para disminuir el tráfico de mensajes a hacia la estación base.

\begin{center}
\begin{figure}[h!]
\centering
 \includegraphics[scale=0.5]{imgs/EstructuraLeach.eps} 
\caption{Estructura de clusters de LEACH.}
\end{figure}
\end{center}

%Descripción detallada

%Descripción del modelo energético
%\subsubsection{Descripción del modelo energético utilizado para el diseño de LEACH}
%LEACH fue ideado para realizar un uso eficiente de la energía basándose en el ``modelo energético de primer orden``. Este modelo plantea que la %energía utilizada para la transmisión de un paquete es la suma entre la energía utilizada en la electrónica de transmisión y la energía utilizada para amplificar la señal para obtener una relación de señal ruido aceptable.
%\begin{center}
% $
%  E_{Tx}(k,d)=E_{elec}*k+\epsilon_{amp}*k*d^2
% $
%\end{center}

%Para la energía utilizada en la recepción, el modelo solo considera la energía utilizada en la electrónica de recepción.
%\begin{center}
% $
%   E_{Rx}(k) = E_{elec}*k
% $
%\end{center}

%Descripción del funcionamiento
\subsubsection{Descripción del funcionamiento}
El algoritmo se construye bajo los siguientes supuestos:
\begin{enumerate}
 \item Cada nodo de la red puede modificar su potencia de transmisión.
 \item Cada nodo de la red tiene comunicación directa con el recolector de datos.
 \item La red es homogénea, todos los nodos tienen los mismos recursos energéticos y de radio transmisión.
\end{enumerate}

LEACH organiza la red en forma jerárquica, dividiéndola en clusters. Cada cluster posee un ``cluster head'' encargado de:

\begin{itemize}
 \item Recolectar los datos entregados periódicamente por los cluster members.
 \item Utilizar una función de agregación para eliminar los datos redundantes.
 \item Enviar los datos agregados hacia el punto de recolección utilizando comunicación directa.
 \item Vía broadcast, proveer a los cluster members una programación TDMA (Time Division Multiple Access), asignándole a cada cluster member una ventana de tiempo para transmitir información.
\end{itemize}

La organización de la red se realiza en forma periódica mediante rondas, cambiando en cada ronda la distribución y composición de los clusters. Cada ronda está dividida en dos fases: Fase Setup y Fase Steady. La fase setup consiste en dos etapas: primero, se realiza la selección del cluster head; segundo, se realiza la formación del cluster. La fase steady consiste en la recolección, agregación y envío hacia el punto de recolección de los datos.\\

El proceso de selección del cluster head asegura que este rol sea asignado  a un nodo durante el tiempo de duración de una ronda y pase a otro cuando comience la siguiente ronda, de manera que el consumo energético asociado a este rol se distribuya por todos los nodos de la red.\\

El proceso de selección del cluster head sigue el siguiente algoritmo.

\begin{enumerate}
 \item Un nodo n cualquiera genera un número flotante aleatorio, v, entre 0 y 1.
 \item El nodo compara v con el ``umbral de selección de cabeza de cluster'',T(n).
 \item El nodo se convierte en cabeza de cluster si v\textless T(n).
\end{enumerate}
 
El valor del umbral de selección de cabeza de cluster, T(n), está diseñado para asegurar que exista una cantidad de nodos, con probabilidad P, que sea elegida como cluster head en una ronda.

\[T(n) = \left\{ 
\begin{array}{l l}
  0 & \quad \mbox{si $n$ }\notin G\\
  P/(1-P*(r mod(1/P))) & \quad \mbox{if $n$ is odd}\\
\end{array} \right. \] 

G representa el conjunto de nodos que ha sido seleccionado como cluster head en las últimas 1 sobre P rondas, y r denota la ronda actual. El parámetro P denota la probabilidad requerida de cluster heads dentro del total de nodos de la red. En general, P es un parámetro arbitrario, sin embargo, existen proposiciones para calcular el parámetro P para optimizar el consumo energético \cite{LeachOptimo}. Cuando el proceso de selección ha terminado, cada cluster head informa su nuevo rol en la red vía broadcast. Cuando un nodo que no ha sido elegido como cluster head recibe el aviso de una cluster head, el nodo debe asociarse a ese cluster, convirtiéndose en cluster member. Si el nodo recibe varios avisos de cluster head, el criterio de decisión es asociarse al cluster head del cuál recibió el aviso con el mayor RSSI (Receive Signal Strength Indication, que indica la intensidad de la señal en dBm) o LQI (Link Quality Indicator, que indica la calidad del enlace de capa MAC, con valor adimensional entre 0 y 255).\\

El protocolo LEACH, en su forma original, considera que todos los nodos tienen comunicación directa con el punto de recolección, por lo que la red no utilizaría una ''comunicación multi-hop real''. Para solucionar este problema se han propuesto algunas modificaciones como: Multi-path, Multi-hop Hierarchical Routing\cite{mumhr}, que es una extensión de LEACH que utiliza comunicación entre cabezas de clusters; y M-LEACH{mleach}, donde se permite que sólo los nodos que tienen comunicación con el punto de recolección tomen el rol de cabeza de cluster, y se plantea un esquema de ''clusters multi-hop'' para el resto de los nodos.

%Protocolo PEGASIS
\subsection{Power-Efficient Gathering in Sensor Information Systems}
%¿Qué es PEGASIS?
Power-Efficient Gathering in Sensor Information Systems (PEGASIS) es un protocolo diseñado para: Extender la autonomía de la red aumentando la eficiencia de la energía y distribuyendo de manera uniforme el consumo a través de todos los nodos de la red; reducir el retraso en que incurre el envío de los mensajes en el camino hacia el punto recolector de datos.\\

%¿Cómo funciona?
Este algoritmo organiza la red en cadenas que realizan el ruteo y agregación de los mensajes que son enviados por los nodos miembros de una cadena hacia el nodo recolector de datos. Dentro de una cadena existen un líder de cadena y el resto funcionan sólo como miembros de la cadena. El líder se encarga de la organización de la cadena y de entregar los mensajes hacia el recolector de datos, mientras que los miembros se encargan de reenviar los mensajes en dirección hacia el líder y realizar agregación de datos para reducir el tráfico. PEGASIS utiliza una rotación circular (estilo round-robin) para distribuir el rol del líder y, en consecuencia, el consumo energético asociado a ese rol. PEGASIS utiliza el mismo modelo energético que se describió en la sección 2.1.1 para LEACH.

\begin{center}
\begin{figure}[h!]
\centering
 \includegraphics[scale=0.5]{imgs/EstructuraPegasis.eps} 
\caption{Estructura de cadenas de PEGASIS.}
\end{figure}
\end{center}

\subsubsection{Descripción del funcionamiento}
El algoritmo de operación funciona de la siguiente manera:
\begin{enumerate}
 \item Un nodo que se considere como el más lejano de el punto de recolección, inicia la construcción de una cadena asociando a su nodo más cercano en dirección hacia el punto de recolección.
 \item El nodo recién asociado, asocia a la cadena al nodo más cercano en dirección hacia el punto de recolección.
 \item El paso 2 continúa hasta que todos los nodos de la red se encuentren asociados a una cadena.
 \item Cuando todos los nodos están asociados, comienza la etapa de rondas. 
\end{enumerate}

Al inicio de cada ronda, se elige el nuevo líder de la cadena utilizando la función de elección de líder, L(r).

\begin{center}
 $
   L_{r}(r) = r \, mod\,  N 
 $
\end{center}
Donde $L_{r}$ corresponde al ID del nodo líder, $N$ al número total de nodos en la red y $r$ a la ronda actual.\\

Esta versión de PEGASIS no cumple el objetivo de disminuir el retraso en la entrega del mensaje hacia el recolector de datos, es más, lo incrementa, por lo que se ha desarrollado otra versión llamada Hierarchical Pegasis\cite{hpegasis}. Además, Hierarchical Pegasis organiza las cadenas por niveles de profundidad y utiliza comunicación entre cadenas para obtener una comunicación multi-hop.

\subsection{Threshold sensitive Energy Efficient sensor Network protocol}
TEEN es un algoritmo jerárquico que recoge las ideas de LEACH (ver sección 2.1.1), división en cluster y comunicación directa aleatoria, llevándolas a un contexto de red reactiva. Las redes de sensores reactivas son aquellas en que un nodo reacciona inmediatamente a cambios repentinos y drásticos del valor del atributo sensado. 

\subsubsection{Descripción del funcionamiento}
El algoritmo de operación funciona de la siguiente manera:
\begin{enumerate}
 \item Organiza la red en clusters, utilizando el algoritmo LEACH.
 \item Periódicamente se distribuye a los cabeza de cluster un umbral para los datos ruteados.
 \item Los miembros del cluster toman mediciones periódicamente y las envían a la cabeza del cluster.
 \item Si el valor medido supera el umbral establecido, la cabeza del cluster reenviará el mensaje hacia el punto de recolección.
 \item El rol de cabeza de cluster rota a través de los miembros del cluster al igual que en LEACH. 
\end{enumerate}

La utilización de un umbral para las mediciones generadas por los nodos no tiene un impacto en la selección de ruta para los mensajes, sino que es utilizada para reducir la magnitud del tráfico hacia el punto de recolección.
%Protocolo HEED
%\subsection{Hybrid, Energy-Efficient, Distributed}
%Hybrid, Energy-Efficient, Distributed (HEED) es un protocolo de clusters distribuídos (similar a LEACH) que provee un algoritmo de control de topología que estructura la red en base a un costo híbrido función de la energía y la comunicación. Los principales objetivos son: Prolongar el tiempo de vida de la red distribuyendo el consumo energía, terminar el proceso de formación de cluster en un número constante de iteraciones, acotar el overhead de control linealmente en función del número de nodos y producir cluster heads bien distribuídas. Con el objetivo de alcanzar  mayor área de cobertura de la red, este algoritmo permite la comunicación inter-cluster.\\

%El diseño de HEED está basado sobre los siguientes supuestos:
%\begin{itemize}
 %\item Todos los nodos tienen control de su potencia de transmisión.
 %\item Los sensores se mantienen en posiciones fijas.
 %\item Los links son simétricos.
 %\item Los nodos no tienen información sobre su posición.
%\end{itemize}

 %En este paper se asume que la principal causa de fallas es el agotamiento de las baterías. JEJE..buena justificación.


%\subsection{Descripción del funcionamiento}
%La primera etapa en la operación de este protocolo es la formación de los clusters.  Para formar un cluster, el protocolo considera los siguientes parámentros.

%\begin{description}
 %\item \textbf{Energía residual:} Puede ser obtenida directamente de un sensor de voltaje o estimada si es que se conocen parámetros como energía consumida por bit sensado, energía para procesamiento, y energía para comunicación.

 %\item \textbf{Costo de comunicación:} Se refiere a un indicador que puede ser función de la proximidad electromagnética entre vecinos o la densidad de un cluster. 
 %\end{description}

%El primer parámetro se utiliza para elegir inicialmente el set de clusters, el segundo se utiliza como variable de decisión para que un nodo cluster member seleccione a que cluster se unirá.

\section{Protocolos de ruteo implementados en TinyOS}

TinyOS provee, por defecto, una pila de comunicación por radio que implementa sólo hasta la comunicación directa. Sin embargo, los desarrolladores unificaron en dos protocolos de red las necesidades más recurrentes de los usuarios\cite{What2ndGenHolds}. Estos protocolos se proveen de forma opcional y el programador debe habilitarlos por medio del cableado de los componentes e instrucciones adicionales en el Makefile.

\subsection{Dissemination }
%¿Qué es diseminación?
Dissemination\cite{TEP118} es un protocolo orientado a mantener la consistencia eventual de variables compartidas a lo largo de la red. Cada nodo mantiene copias locales de las variables compartidas, cuando el servicio de diseminación le informa a un nodo que el valor se ha actualizado, este intercambia paquetes con un nodo que tenga el valor actual. Así, luego de un periodo de tiempo, todos los nodos tendrán el valor actualizado.\\

%¿Cómo funciona?
La implementación de Dissemination utiliza el algoritmo Trickle \cite{TricklePaper} para comunicar los valores de las variables compartidas. La idea de Trickle es que periódicamente los nodos comunican vía broadcast su copia del set de variables compartidas a sus vecinos, pero no realizan esta función si es que han recibido recientemente una copia idéntica de las variables compartidas a las que tienen almacenada. Cuando un nodo recibe una copia antigua del set de variables compartidas, realiza un broadcast para actualizar el set de sus vecinos. De esta forma, los nodos mantienen actualizadas las variables compartidas utilizando una inundación controlada de los mensajes de actualización.

%Carácterísticas
%Me aburrió también


%\subsection{TYMO}
%¿Qué es TYMO?
%Tymo es la implementación del protocolo DYMO (ref) sobre TinyOS. Es un protocolo punto a punto (aclaración), que mantiene pequeñas tablas de ruteo para cumplir con las limitaciones de memoria de un dispositivo embebido.\\ 
%¿Cómo funciona?
%Cuando un nodo necesita una ruta, disemina un mensaje Route Request(RREQ), el mensaje se distribuye por la red vía flooding(aclaración). Cuando el nodo llega al nodo destino o un nodo que tenga una ruta actualizada hacia el destino, responde con un Route Reply (RREP). TYMO permite que un mensaje RREQ o RREP acumule información de la ruta que sigue dentro del mismo mensaje, así un nodo origen, destino o intermedio puede obtener una gran cantidad de información de ruteo cuando recibe el mensaje. Cuando las rutas no son usadas por mucho tiempo, son eliminadas para que las entradas relacionadas en la tabla puedan ser reutilizadas. DYMO utiliza variable simples para determinar la calidad de una ruta, como el número de secuencia de los mensajes y la cuenta de saltos.

%Se puede rellenar con un poquito más cierto....

\subsection{Collection Tree Protocol}
%¿Qué es CTP?
CTP \cite{TEP123} provee una comunicación del tipo ''mejor esfuerzo'' entre un nodo y alguna raíz en la red, con la cual el nodo está relacionado. Es un protocolo que utiliza topología de árbol: algunos nodos pueden asumir el rol de raíz, recolectando los datos sensados por todos sus nodos hijos; los nodos que conforman las ramas se ocupan de tomar mediciones y reenviar hacia sus padres la información recolectada por sus hijos.\\

%¿Cómo funciona?
En CTP se forma un árbol de recolección, mediante asociaciones que utilizan el RSSI como variable de selección. Es un protocolo sin direcciones, ya que todos los mensajes son direccionados a la raíz. Utiliza tablas de ruteo que almacenan un indicador acumulativo que evalúa la congestión de una ruta (el ETX) como gradiente para establecer el próximo salto de un mensaje.\\

Este protocolo ha sido especificado en el TEP 123 (TinyOS Enhancement Proposals) para su implementación. La especificación se ha finalizado en Enero del año 2009 y actualmente se están desarrollando pruebas sobre la implementación que ya está incluida en la distribución de TinyOS 2.X.

%\subsection{Berkeley Low-Power IP Stack}
%Damn it, este protocolo es bakan y cumple todo.
%¿Qué es BLIP?
%Berkeley Low-Power IP Stack (BLIP) es un protocolo que implementa IPv6 para redes de dispositivos TinyOS. Actualmente permite formar redes IP multihop con nodos que utilizan TinyOS y permite interoperabilidad con redes con otras redes IP, sin embargo, no cumple el estándar completamente. BLIP es la segunda implementación de este protocolo sobre TinyOS 2.x y en la última versión ha incluído soporte para el servicio LPL de TinyOS (aclaración).\\

%¿Cómo funciona?
%Es una red IP existen dispositivos hosts y routers. Cada mote que corre el stack BLIP funciona como un router, es capáz de reenviar paquetes y tomar decisiones de ruteo. Además, en BLIP existen los ``Edge routers'', que son los routers que pueden comunicar una red con otra, es decir, estos routers pueden funcionar como gateway para la red de sensores inalámbricos.

%Agregar características ---ya me aburrí.

\section{Selección de un esquema de recolección}
%De lo general a lo particular
Luego de investigar los algoritmos el objetivo es seleccionar el esquema de recolección que pueda generar los mejores resultados para la aplicación  descrita en la sección 1.3. Dado que la principal variable de diseño para un protocolo de redes de sensores inalámbricos es la energía consumida por los nodos y por la red, las métricas para comparar los algoritmos serán las relacionadas con el consumo energético. Se excluye el algoritmo TEEN, dado que es una modificación mínima a LEACH y que la aplicación para el sistema de monitoreo requiere que las mediciones sean recolectadas periódicamente.

%Sección para comparar protocolos propuestos e implementados sobre sus métricas
\subsection{Análisis comparativo de resultados}
Para comparar los algoritmos descritos en la sección 2.1, se revisan publicaciones para determinar cuáles son las principales métricas para evaluar los protocolos descritos. La principales métricas que se utilizan para evaluar algoritmos de ruteo en la literatura son:

\begin{itemize}
\item \textbf{Porcentaje de nodos muertos, según la cantidad de rondas:} Esta métrica es la que se encuentra en todas las publicaciones que analizan estos protocolos. Como se ha descrito antes, una ronda corresponde al intervalo regular de tiempo en el cual la topología de la red se mantiene estática. Por otro lado, el concepto de ''nodo muerto'' varía según la publicación pero, en general, se refiere a un nodo que ha utilizado tanta cantidad de energía, que ya no puede realizar los procesos de comunicación inalámbrica.

\item \textbf{Delay (retardo):} Esta métrica se refiere a la cantidad de tiempo que transcurre entre la generación de un paquete de datos en un punto específico de la red y el arribo de ese mensaje al punto de recolección de la red.

%\item \textbf{Energía total disipada en la red:} Esta métrica se refiere a la energía disipada por la red, obtenida mediante la utilización de modelos energéticos, que proveen un valor para la energía en función de: distancia entre nodos, largo de los mensajes, supuestas características eléctricas y condiciones energéticas iniciales.
\end{itemize}

Según los resultados, los retardos en la transmisión de un mensaje desde fuente a destino son del orden de los milisegundos, en consecuencia, no son significativos frente a la tasa de generación de mensajes que se utilizan, por ejemplo, en el monitoreo de variables de suelo para aplicaciones relacionadas con la agronomía, donde la tasa de generación de mensajes es del orden de un mensaje por minuto. Por esta razón, se privilegian los resultados de las métrica ''porcentaje de nodos muertos, según la cantidad de rondas''.\\

En la publicación original que introduce el algoritmo PEGASiS, se realiza una comparación entre el algoritmo LEACH y PEGASIS. Los resultados que se muestran en la figura 2.3 fueron obtenidos realizando una simulación en el simulador de redes NS2, utilizando una energía inicial de 1 [J] y un área de despliegue de 100[$m^2$].

\begin{center}
\begin{figure}[H]
\centering
\includegraphics[scale=0.3]{imgs/AnalisisComparativoProtocolosPAPERPEGASIS.eps}
\caption{Comparación entre LEACH y PEGASIS de la publicación \cite{PegasysOriginal}, utilizando el número de rondas en función del porcentaje de nodos muertos.}
\end{figure}
\end{center}

Según los resultados de esta publicación, se concluye que PEGASIS  permite un tiempo de operación de los nodos más extenso que LEACH, e incluso lo triplica.\\

En la publicación \cite{PaperComparacion} se realiza una evaluación de los algoritmos LEACH y PEGASIS. Para esta evaluación se utiliza el simulador SENSE, desplegando 2000 nodos en un área de 100[$m^2$], con una tasa de generación de mensajes de 1 mensaje por segundo y una energía inicial de 2[J].

\begin{center}
\begin{figure}[H]
\centering
\includegraphics[scale=0.3]{imgs/ComparativaRondasVsMuertos.eps} 
\caption{Comparación entre el número de rondas y el porcentaje de nodos muertos}
\end{figure}
\end{center}
%También, en términos de energía disipada en la red, PEGASIS es superior a LEACH, como muestra la figura 2.5.
%\begin{center}
%\begin{figure}[H]
%\centering
%\includegraphics[scale=0.375]{imgs/ComparativaEnergiaDisipada.eps} 
%\caption{Comparación entre la energía total disipada en la red por los algoritmos}
%\end{figure}
%\end{center}
Como se muestra en la figura 2.4, el algoritmo PEGASIS es el que permite un mayor número de ronda, con más nodos vivos en la red. Con estos resultados, se decide seleccionar el algoritmo PEGASIS para la implementación del protocolo de red. La implementación de este protocolo se describirá en el capítulo 3, sin embargo, dado que fue insuficiente en términos de generar una red multi-hop y desempeño energético, en el capítulo 4 se describe el diseño de un protocolo de árbol de recolección que permite satisfacer las metas propuestas en este trabajo.

%\section{Dificultades de esta etapa}

	
	%Para armar Indice que quiere el profe
	%archivo_indice.tex
\chapter{Anexos}
	\section{Hardware}
		\subsection{Macintosh}

		\subsection{iPhone, iPad, iPod Touch}
		descripcion de iphone, cuando salió, presentación de steve, revolucion del smartphone, pasar de teclas a nuevas pantallas táctiles.
		ipad, nueva forma de leer, reemplaza netbook, 3g pero no llama, creó un nuevo mercado y es lider
		ipod, comezó como música, luego de iphone adopta ios, explicar que es alternativa a iphone, gran adopcion en usa.  lider en players.
	\section{Sistemas operativos}
		\subsection{OSX}
		os x, se empezó el desarrollo en lion, migración de xcode 3 a 4, inestabilidad
		explicar como se usa  itunes con ios
		\subsection{iOS}
		se empezó con ios 4, pasamos a 5, luego 6		
		facilidad de actualizacion
	\section{Frameworks}
	uikit, descripción 
	avfoundation, descripcion
	mpmedia framework, descripcion
	%entorno de desarrollo utilizado
	
%\chapter{Bibliografia}
	% refencias de lo estudiado y utilizado para el trabajo	
	
	%Incluir capítulo 3
	%	
\chapter{Pruebas de funcionamiento}

Para revisar el buen funcionamiento de la aplicación desarrollada, se utilizaró el software dispuesto por Apple Inc,  el entorno de desarrollo Xcode y herramientas asociadas. Para obtener estas aplicaciones, es necesario poseer una cuenta de desarrollador para \textbf{iOS} o \textbf{Mac OS X}, las cuales se pueden obtener en el sitio web \url{https://developer.apple.com/programs/} luego de pagar la matrícula impuesta.\\

La cuenta de desarrollador permite acceso al repositorio de herramientas \cite{apple-repositorio} y también acceso a versiones de los sistemas operativos antes de lanzamiento.
Además permite poner a la venta las aplicaciones desarrolladas en la \textit{\textbf{App Store}}\cite{apple-appstore}.

% datos en el camino, seguimiento de un tweet, pantallazos de debug
\section{Cambio de tiempo y/o canal}

Las primeras pruebas estuvieron relacionadas con el cambio de tiempo a través de peticiones al dispatcher con argumentos en la cadena de consulta, estas se realizaron con la aplicación de OS X: \textbf{Quicktime X}, la cual permite reproducir streams que cumplen con las espeficicaciones de HTTP Live Streaming. Otra alternativa que además se encuentra disponible en otros sistemas operativos es \textbf{VideoLan VLC}. Las pruebas consistieron en modificar los argumentos \textbf{t}, \textbf{s} y \textbf{c} de URLs del tipo:
\url{http://ssdemo.altavoz.net/playlist/playlist.m3u8?s=tvn&t=1353186024&c=0}, donde el parámetro \textbf{t} corresponde al tiempo Unix del 17 de noviembre de 2012 18:00:24 GMT-3.\\

El escificar el tiempo en el parámetro \textbf{t} resulta siempre en una lista de reproducción con segmentos asociados al tiempo correspondiente. Además se monitoreó la respuesta a través de la aplicación rastreadora de paquetes Wireshark. Para esto se necesitó aplicar un filtro especial de forma que los paquetes relacionados con HTTP Live Streaming fueran entregados. El filtro corresponde a mostrar paquetes recibidos desde el servidor donde el cliente, en este caso Quicktime, se contacta para obtener una lista de reproducción, y además pedir la entrega de paquetes que tengan en su contenido la palabra \textbf{apple}.

\begin{figure}[H]
	\centering
\begin{lstlisting}
		ip.addr==200.91.44.49 && http contains apple
\end{lstlisting}
\caption{Filtro en Wireshark que muestra sólo paquetes relacionados con el HLS}
\end{figure}


Los resultados obtenidos son los esperados y corresponden a la reproducción del video en el momento indicado y en el canal indicado según los parámetros.

\begin{figure}[H]
	\centering
	\begin{tabular}{cc}
	\includegraphics[scale=0.3]{imgs/qt-mega.png} & 
	\includegraphics[scale=0.3]{imgs/qt-tvn.png} \\
	\end{tabular}
	\caption{Captura de Quicktime reproduciendo lista de reproducción a través de un URL.}
	\label{qt-tvn-mega}
\end{figure}

Utilizando \textbf{Wireshark} se observaron los datos entregados al cliente desde el servidor como \textit{\textbf{cookies}} indicadas en \ref{subsec:cookies}.

\begin{figure}[H]
	\centering
\begin{lstlisting}
HTTP/1.1 200 OK
Server: nginx/1.2.0
Date: Sat, 24 Nov 2012 18:06:10 GMT
Content-Type: application/vnd.apple.mpegurl
Transfer-Encoding: chunked
Connection: keep-alive
Keep-Alive: timeout=20
X-Powered-By: PHP/5.3.15-1~dotdeb.0
Set-Cookie: stream=tvn
Set-Cookie: t0=1353186024
Set-Cookie: tp=1353780370
Set-Cookie: ti=1353186033
Set-Cookie: tl=1353186085
Set-Cookie: si=1
Set-Cookie: sf=6

#EXTM3U
#EXT-X-TARGETDURATION:10
#EXT-X-PROGRAM-DATE-TIME:2012-11-17T18:00:33-03:00
#EXT-X-VERSION:3
#EXT-X-MEDIA-SEQUENCE:1
#EXTINF:13,
http://ssdemo.altavoz.net/streams/tvn/tvn_1353186033.ts
#EXTINF:10,
http://ssdemo.altavoz.net/streams/tvn/tvn_1353186043.ts
#EXTINF:9,
http://ssdemo.altavoz.net/streams/tvn/tvn_1353186052.ts
#EXTINF:12,
http://ssdemo.altavoz.net/streams/tvn/tvn_1353186063.ts
#EXTINF:7,
http://ssdemo.altavoz.net/streams/tvn/tvn_1353186070.ts
#EXTINF:13,
http://ssdemo.altavoz.net/streams/tvn/tvn_1353186085.ts
\end{lstlisting}
\caption{Información HTTP entregada por servidor Web}
\label{lst:setcookies}
\end{figure}

De la figura \ref{lst:setcookies} se puede observar que el servidor web Ngnix\ref{bib:ngnix-homepage} entrega en la cabecera HTTP las intrucciones Set-Cookie con los valores descritos en \ref{subsec:cookies}. \\

En la siguiente conexión con el servidor, se ajustaran los valores de las \textit{cookies} acorde al contenido de la lista de reproducción entregada. Como se describió anteriormente las \textit{cookies} no son manejados por la aplicación cliente, ya que el sistema operativo iOS se encarga de manejar los valores de estas.\\

En la figura \ref{lst:sequence2} se puede observar que los nuevos valores de cookies son ajustados por el servidor para la lista de reproducción acorde a la segunda secuencia del stream.

\begin{figure}[H]
	\centering
\begin{lstlisting}
HTTP/1.1 200 OK
Server: nginx/1.2.0
Date: Sat, 24 Nov 2012 18:06:25 GMT
Content-Type: application/vnd.apple.mpegurl
Transfer-Encoding: chunked
Connection: keep-alive
Keep-Alive: timeout=20
X-Powered-By: PHP/5.3.15-1~dotdeb.0
Set-Cookie: stream=tvn
Set-Cookie: tp=1353780385
Set-Cookie: ti=1353186043
Set-Cookie: tl=1353186093
Set-Cookie: si=2
Set-Cookie: sf=7

#EXTM3U
#EXT-X-TARGETDURATION:10
#EXT-X-PROGRAM-DATE-TIME:2012-11-17T18:00:43-03:00
#EXT-X-VERSION:3
#EXT-X-MEDIA-SEQUENCE:2
(...)
\end{lstlisting}
\caption{Información HTTP entregada por servidor Web con la lista actualizada para secuencia 2}
\label{lst:sequence2}
\end{figure}



%http://ssdemo.altavoz.net/playlist/playlist.m3u8?s=tvn&t=1353186024&c=0
%http://ssdemo.altavoz.net/playlist/playlist.m3u8?s=video1&t=1353186024&c=0

%escribir en consola el URL, probarlo en quicktime, vlc, navegador, app, wireshark para ver cookies

\section{Proceso corriendo en fondo} % background

Debido a que se utiliza el protocolo HTTP Live Streaming de forma que entregue información en vivo, la instancia de AVPlayer se mantiene realizando peticiones por nuevas listas de reproducción a medida que pase el tiempo, sin importar que la reproducción se encuentre \textbf{pausada} o que la aplicación esté corriendo en el \textbf{fondo} (\textit{background}).\\

Se encontró un problema al interrumpir la reproducción. Sea cual sea el motivo, por ejemplo recibir una llamada telefónica, el sistema iOS envía la aplicación cliente a background pausando la reproducción. Al volver despues de la interrupción la aplicación retoma la reproducción, sin embargo el tiempo que ha pasado implica en  listas de reproducción distintas a la que se estaba viendo.\\

La solución fue establecer una marca de tiempo de 30 segundos, acordes a la ventana de transmisión de la lista, asociada a la instrucción de pausa y reproducción del reproductor, que además es llamada por el botón \textit{Play/Pause} en la interfaz gráfica.\\

Al enviar la aplicación a background se guarda la \textbf{fecha del contenido} del stream y el momento en el cual se ha pausado.
Al volver a \textit{foreground} se activa el método para reanudar la reproducción donde se compara el momento de pausa guardado y el actual. Si esta diferencia es mayor a 30 segundos, se realiza una nueva llamada al servidor mediante URL. Sin embargo cuando la diferencia es menor, se reanuda la reproducción debido a que la fecha del contenido del stream aun se encuentra dentro del rango en la lista que posee la instancia de AVPlayer.\\

El valor de 30 segundos se debe una lista de reproducción con mínimo 3 segmentos de 10 segundos.\\

El problema encontrado se obtuvo gracias a la depuración de la aplicación ejecutandose en un dispositivo iOS (iPhone). Se revisó el valor del \textbf{rango de reproducción} a través de notificaciones por cambios (página \pageref{item:seekableTimeRanges}). A pesar que la reproducción se pausaba o la aplicación se mandaba a background el rango se mantenia en aumento.

%escribir en consola los seekable time ranges, poner pausa background y al retomar reproducir, se notó que avanzaba la lista de iguañl forma. se arregó guardando la fecha al momento de poner pausa. que se llamaba al irse a background.

\section{Cambio de ancho de banda}

El cambio de ancho de banda se revisó mediante Wireshark y modificando el valor de la cadena de consulta \textbf{c} entre 0 (WiFi) y 1 (red celular) en el URL.\\ 

Por ejemplo: \url{http://ssdemo.altavoz.net/playlist/playlist.m3u8?s=tvn&t=1353186024&c=0}
La respuesta fue rastreada obteniendo una lista de reproducción:

\begin{figure}[H]
	\centering
\begin{lstlisting}
#EXTM3U
#EXT-X-STREAM-INF:PROGRAM-ID=1,BANDWIDTH=320000
stream.m3u8?s=tvn&t=1353186024
#EXT-X-STREAM-INF:PROGRAM-ID=1,BANDWIDTH=65000
stream.m3u8?s=tvn&a=1&t=1353186024
\end{lstlisting}
\caption{Lista de reproducción resultante de una petición al servidor con parámetro \textbf{c = 0}}
\label{lst:playlistc0}
\end{figure}

\begin{figure}[H]
	\centering
\begin{lstlisting}
#EXTM3U
#EXT-X-STREAM-INF:PROGRAM-ID=1,BANDWIDTH=65000
stream.m3u8?s=tvn&a=1&t=1353186024
#EXT-X-STREAM-INF:PROGRAM-ID=1,BANDWIDTH=320000
stream.m3u8?s=tvn&t=1353186024
\end{lstlisting}
\caption{Lista de reproducción resultante de una petición al servidor con parámetro \textbf{c = 1}}
\label{lst:playlistc1}
\end{figure}

En el caso de ajustar el valor a 1, la lista de reproducción invierte la prioridad de las variantes del stream, priorizando el enlace que posee sólo audio para un ancho de banda de \textbf{65kbps}. \ref{lst:playlistc1}

%se especifica requerimiento a la empresa, se modifica el dispatcher con lista de variantes, se revisa en wireshark.
%luego se utiliza la herramienta network link conditioner para modificar el BW en plena transmisión, revisar wireshark cuando cambia y ver en la misma app.
%Luego se prueba con iPhone cambiando entre wifi y 3g

  \subsection{WiFi}
  se agrega campo en el URL para priorizar distintas variantes, con wifi video, 3g audio.
  explicar que parte con el video, mostrar 2 imagenes emparejadas.
  \subsection{3G}
  con 3g, imagen de que muestra el audio primero.
\section{Cambio de transmisión mediante Twitter}
- se revisa usando safari desktop con script que incluye player

  \subsection{Dentro de la aplicación}
	- debug viendo el cambio de tranmisión, URL generado y definiendo el mismo canal. manejo en caso de tweet sin enlace
	se genera notificación interna con los componentes. ver que cambia el stream.
  \subsection{Scheme registrado en iOS}
	- clientes de twitter en iPhone, twitter y Tweetbot, debug en el método e imprimir en consola la app que lo lanzó.  
	- debug con app en background
	- debug con app cerrada y a la espera que parta por otra aplicación
	explicar que se detuvo la carga del stream cuando está cerrada y que se espera una notificacion de la carga de la UI antes de cargar el stream.
  
\section{Enlaces Twitter en otros dispositivos}
	pantallazos de páginas de error, el por qué y cómo se entregan según el script, no se revisa internet explorer porque se utilizaron navegadores compatibles con OS X.
	
  \subsection{PC Escritorio}
    \subsubsection{Chrome}
    chrome no es compatible con hls
    foto
    \subsubsection{Safari}
    safari muestra un player html5
    foto del player
    \subsubsection{Firefox}
    firefox no es compatible
    foto
    \subsubsection{Opera}
    opera tampoco
    foto
  \subsection{Android y otros móviles incompatibles}
  pantallazo del cell de la esperanza
  explicar que tampoco es compatible con HLS, por lo menos 2.3, 4.0 ICS en adelante es compatible, pero para reproducir el stream se debe desarrollar una aplicación cliente similar a la hecha en iOS
\section{Comparación con SocialStream Flash}
mostrar player original de eduardo, la diferencia de la linea de tiempo, el salto de tiempo, la exactitud y que no es tan exacto porque los segmentos son de 10 segundos

la linea de tiempo difiere con una representación gráfica, se utilizó el control de fecha para simplificar el punto a elegir, además la linea de tiempo precisa de mouse over para mostrar datos y en el caso táctil difiere mucho.



	
	%Incluir capítulo 4
	%	\chapter{Diseño de un protocolo con topología de árbol}

Debido a los resultados obtenidos en la evaluación energética de la implementación del protocolo PEGASIS descrita en el capítulo 3, se decidió diseñar e implementar un nuevo enfoque de recolección de datos, utilizando una técnica que permita reducir el consumo del TMoteSky en los tiempos de idle listening. Se optó por implementar un árbol de ruteo, similar a la construcción de red propuesta en el estándar IEEE 802.15.4, construido para utilizar la herramienta Low Power Listening para reducir el consumo de corriente del dispositivo. \\

\section{Visión general del diseño}

El diseño comienza con la búsqueda de una estrategia de ahorro energético que permita reducir el consumo de corriente del TMoteSky, atacando el problema del consumo de corriente del chip CC2420 en estado idle. Posteriormente, se diseña el '' Componente de recolección'', un esquema de formación y control de la topología, que permite establecer una red con topología de árbol multihop y mantener la recolección de datos incluso si se presentan nodos muertos. 

\section{Estrategias de ahorro energético}
El perfil del consumo de corriente del TMoteSky, está dominado principalmente por el consumo del transceptor CC2420.

\begin{table}[H]
  \centering
  \caption{Principal consumo de corriente en el TMoteSky  \cite{TelosBConsumo}}
 \includegraphics[scale=0.5]{imgs/ConsumoTmoteSky.eps}
\end{table}

Como se muestra en el cuadro 4.1, el consumo del CC2420 representa más del 90\% del consumo total del dispositivo.\\

Para los dispositivos que soportan TinyOS, se han planteado dos estrategias de ahorro energético: La primera consiste en utilizar un ciclo de trabajo sincronizado, donde los dispositivos mantienen sus radios apagadas por periodos largos de tiempo y las enciendan por periodos cortos, realizando este ciclo de manera sincronizada en toda la red; o, utilizar la técnica ''Low Power Listening'', que consiste en apagar la radio y encenderla periódicamente, por periodos de tiempo cortos, para sensar el canal y detectar cuando otro dispositivo quiere transmitir un mensaje.\\

La primera estrategia requiere un protocolo de sincronización de tiempo, que consiste en un conjunto de operaciones y mensajes que se intercambian en la red para mantener secciones o la totalidad de la red sincronizada en el tiempo. TinyOS provee un protocolo llamado ''Flooding Time Sync Protocol'' (Protocolo de sincronización de tiempo por inundación, FTSP), que consiste en inundar la red con el tiempo del nodo con menor TOS\_NODE\_ID en la red (''tiempo global''). Sin embargo, se ha documentado que la implementación para la radio CC2420 hereda un bug de la implementación de la interfaz PacketTimeStamp para el CC2420\cite{FtspBug}. El error causa que continuamente los receptores del mensaje de sincronización obtengan un tiempo global errado por decenas de milisegundos, lo que eventualmente causa la pérdida de sincronización.\\

La estrategia ''Low Power Listening'' tiene el objetivo de disminuir el consumo de corriente, disminuyendo el tiempo de idle listening y a su vez permite que la red sea asincrónica. Para lograr estos objetivos, la implementación toma las medidas siguientes:

\begin{itemize}
\item Se implementa un ciclo de poder para la radio (ciclo LPL), es decir, ciclo de radio encendida y radio apagada.
\item La radio se mantiene encendida el tiempo suficiente para detectar un carrier en el canal.
\item Si detecta un carrier, se mantiene encendida la radio para recibir el mensaje.
\item Cuando se envía un mensaje, se extiende el tiempo de duración del preámbulo hasta el tiempo de duración del ciclo LPL.
\end{itemize}

\begin{center}
\begin{figure}[H]
\centering
\includegraphics[scale=0.625]{imgs/GraficoTiemposLPLsoloLPL.eps} 
\caption{Gráfico del consumo de corriente que tiene el TmoteSky con módulo LPL activado, cuando realiza operaciones escucha en bajo consumo.}
\end{figure}
\end{center}

\begin{center}
\begin{figure}[H]
\centering
\includegraphics[scale=0.625]{imgs/GraficoTiemposLPLTX.eps} 
\caption{Gráfico del consumo de corriente que tiene el TmoteSky con módulo LPL activado, cuando transmite un paquete con extensión de preámbulo.}
\end{figure}
\end{center}

 Debido a la ausencia de un protocolo robusto de sincronización de tiempo implementado para TinyOS y a que implementar un protocolo de esta naturaleza está fuera del alcance de este trabajo, se decidió utilizar la estrategia ''Low Power Listening'' para gestionar el consumo de corriente del TMoteSky. En el  Anexo C, se describe en detalle la utilización de este componente y se realiza una caracterización del consumo para distintos escenarios.

\section{Diseño del componente de recolección}

Al igual que la implementación del protocolo PEGASIS, se construye el protocolo de red sobre la pila de comunicación ''Radio Stack'' que provee TinyOS. Además, se decide incluir otras mejoras no exploradas antes: Una cola de reenvío, para reducir las posibilidades de perder mensajes; y mensajes de datos con carga de largo variable, para dar libertad a la aplicación usuaria del protocolo a que utilice una estructura de datos arbitraria para sus datos. Siguiendo el mismo esquema de implementación que CTP y Dissemination, el protocolo debe implementar las funciones Send y Receive, que proveerán las funciones de red necesarias para que el nodo pueda enviar y recibir mensajes al Gateway.
%Diagrama del Stack
%\begin{center}
%\begin{figure}[H]
 %\centering
 %\includegraphics[scale=0.4]{imgs/Stack.eps}
 %\caption{Stack}
%\end{figure}
%\end{center}
\subsection{Formato de mensajes}
En las próximas secciones se mencionarán de forma recurrente los mensajes definidos para esta implementación, por lo que a continuación se describen los dos tipos de mensajes diseñados para el protocolo árbol: Mensajes de Topología y Mensajes de Datos Multihop.\\ 

Los mensajes de topología son todos los mensajes intercambiados para la formación y mantención de la topología de árbol de la red, a los cuales se les asigna el AM Type 0x20.
%Topology Message
\begin{table}[H]
  \centering
  \caption{Descripción de la estructura de un mensaje de topología.}
 \includegraphics[scale=0.5]{imgs/MensajesProtocoloMultihopTopologia.eps}
\end{table}

Dentro de los mensajes de topología establecen los siguientes tipos: ''beacon request'', ''beacon response'', ''join request'', ''join response ok'' y ''join response denied''.

%Multihop Message Header
\begin{table}[H]
  \centering
  \caption{Descripción de la estructura de un mensaje de datos.}
 \includegraphics[scale=0.5]{imgs/MensajesProtocoloMultihopMensajeDato.eps}
\end{table}

Los mensajes de datos multi-hop son los que transportan los datos generados por la aplicación a través de la red, a los cuales se les asigna el AM Type 0x21. Para estos mensajes sólo se define la cabecera, ya que el resto del mensaje variará según como sea definido por la aplicación.

%Secuencias
\subsection{Descripción de secuencias y lógica}
En esta sección se describen las secuencias de mensajes y la lógica que se implementa para la formación de una red con topología de árbol. Para formar el árbol de asociación se utiliza el esquema planteado por la IEEE 802.15.4 \cite{IEEE802154}, 
%Secuencia de Asociación
\subsubsection{Secuencia de Asociación}
Para formar el árbol de recolección se utiliza un sistema de beacons (faros). Los nodos descubren la red enviando un mensaje de ''Beacon Request'', que sólo será respondido por un nodo que ya se encuentre en la red. El Gateway asume el rol de raíz y primer beacon, y en consecuencia, a partir del descubrimiento del Gateway, la red se ramifica en forma de árbol. La secuencia de asociación que propone este esquema consiste en el intercambio de cuatro mensajes:\\
 
 \textbf{BEACON\textunderscore REQUEST:} Este mensaje es utilizado para permitir que el nodo emisor encuentre la red. Debe incluir al menos la dirección origen, para permitir que los nodos que estén en la red le respondan con un mensaje BEACON\textunderscore RESPONSE, y debe ser direccionado a la dirección de broadcast.\\
 
 \textbf{BEACON\textunderscore RESPONSE:} Este mensaje se envía en respuesta a un BEACON\textunderscore REQUEST. Sólo puede ser enviado por un nodo que ya esté asociado a la red y debe proveer la información necesaria para que el nodo que está ingresando a la red pueda decidir entre varios mensajes BEACON\textunderscore RESPONSE el que estime como el mejor candidato para negociar una asociación.\\
 
 \textbf{JOIN\textunderscore REQUEST:} Si el nodo que está tratando de asociarse a la red ha recibido mensajes BEACON\textunderscore RESPONSE y ha determinado un nodo candidato para realizar la negociación de asociación, envía un mensaje de JOIN\textunderscore REQUEST dirigido al nodo candidato. Este mensaje debe contener al menos la dirección de origen del nodo que realiza este requerimiento.\\
 
 \textbf{JOIN\textunderscore RESPONSE:} Un nodo que está en la red que recibe un JOIN\textunderscore REQUEST debe responder al nodo que ha solicitado la asociación con la red en forma positiva o negativa. Este mensaje debe contener al menos la dirección origen y la respuesta a la asociación.\\
 
 %Diagrama de secuencia de asociación
\begin{figure}[H]
 \centering
 \includegraphics[scale=0.45]{imgs/DiagramaSecuenciaAsociacion.eps}
 \caption{Diagrama UML de secuencia de asociación}
\end{figure}

En la figura 4.3 se muestra la secuencia que se desea implementar para asociar un nodo que entra a la red (Nodo A) con otro que ya está en la red (Nodo B). Primero, nodo A envía un mensaje broadcast BEACON\textunderscore REQUEST, el nodo B escucha el mensaje y envía un BEACON\textunderscore RESPONSE. El nodo A espera un tiempo arbitrario, denominado "BEACON\_REQUEST\_INTERVAL" , que permite que el nodo A reciba y almacene temporalmente varios mensajes BEACON\textunderscore RESPONSE. Si el nodo ha recibido mensajes beacon request y ha determinado un candidato para solicitar la asociación a la red, envía un JOIN\textunderscore REQUEST al candidato. Si no ha encontrado ningún candidato, vuelve a enviar un BEACON\textunderscore REQUEST.

%Secuencia de Multihop
\subsubsection{Secuencia de envío y reenvío de mensajes multihop}
En esta secuencia se intercambian dos tipos de mensajes.\\

\textbf{MH\textunderscore DATA:} Estos mensajes pueden ser generados por un nodo que está asociado a la red y son los que transportan datos entregados por la capa de aplicación. Estos mensajes deben ser capaces de ser generados con largo variable, para permitir a la aplicación que envíe dentro de estos mensajes las estructuras de datos que estime convenientes. La cabecera debe ser de un largo definido y debe al menos incluir: la dirección de origen del mensaje, un número de secuencia y el largo del mensaje. La dupla formada por la dirección del mensaje y número de secuencia funciona como una ''clave primaria'' que permite identificar, de forma única, un mensaje. La identificación de un mensaje es útil para tolerar fallas y la organización de los mensajes en el destino final, por ejemplo, una base datos. Estos mensajes pueden ser generados por un nodo, o bien, recibidos por un nodo para posteriormente ser reenviados.\\

\textbf{ACK:} Este tipo de mensaje debe ser generado cada vez que un nodo recibe un mensaje MH\textunderscore DATA para informar al nodo fuente del mensaje, que el nodo intermedio o destino ha recibido el mensaje correctamente. Estos mensajes son de sólo un salto, ya que como este protocolo es diseñado para realizar el "mejor esfuerzo" no se implementan ACKs multihop.\\

%Diagrama de secuencia multihop
\begin{figure}[H]
 \centering
 \includegraphics[scale=0.45]{imgs/DiagramaSecuenciaMultihop.eps}
 \caption{Diagrama de secuencia de asociación}
\end{figure}

En la figura 4.4 se observa la secuencia de envío de un mensaje Multihop. El nodo A determina que su próximo salto es el nodo B y envía un mensaje MH\textunderscore DATA al nodo B. El nodo B responde con un ACK si recibe el mensaje correctamente, sino no envía mensaje alguno.  Luego el nodo B determina que el próxima salto es el nodo C y reenvía el mensaje recibido hacia el nodo C. Está secuencia debe continuar a través de todos los nodos intermedios hasta que finalmente llega al Gateway.

%Secuencia de Reconexión
\subsubsection{Secuencia de reconexión}
Puede suceder que un nodo presente algún tipo de falla, agote su batería o simplemente desaparezca aparentemente de la red. En ese caso, los nodos que rutean mensaje a través de él, deben detectar que este nodo ya no está presente. Ya que los nodos cuando reciben un mensaje de datos deben confirmar con un ACK, un nodo puede identificar que un nodo ha desaparecido aparentemente por la pérdida de ACKs.\\

%Diagrama de Reconexión
\begin{figure}[H]
 \centering
 \includegraphics[scale=0.45]{imgs/DiagramaSecuenciaReconexion.eps}
 \caption{Diagrama de secuencia de renovación de la asociación}
\end{figure}

La secuencia de desconexión y reconexión comienza con la pérdida sucesiva de ACKs, digamos que el nodo A está ruteando mensajes a través del nodo B, y el nodo B aparentemente desaparece. Cuando el nodo A intenta enviar un mensaje MH\textunderscore DATA al nodo B,éste no responde con un ACK, por lo tanto el nodo A puede identificar un ACK perdido. Luego de un número definido arbitrariamente de ACKs perdidos, el nodo A deja la asociación con el nodo B. Posteriormente, el nodo A vuelve a buscar un padre utilizando la secuencia de asociación descrita anteriormente.

%Secuencia BROADCAST
\subsubsection{Secuencia de MHData Broadcast}
 Es posible, que la aplicación requiera enviar un mensaje de datos para configurar la aplicación de uno, varios o todos los nodos. Al no utilizar direcciones, ni tablas de ruteo que indiquen el camino a un nodo en particular, no se puede generar una ruta específica a un nodo X desde el Gateway. Otros protocolos similares de recolección (ver CTP más Dissemination) realizan esta funcionalidad mediante una inundación o diseminación del mensaje, a través de toda la red. Para emular este comportamiento, se decide diseñar una funcionalidad de "broadcast de red" para los mensajes de datos, que pueda ser utilizada por cualquier nodo de la red. El esquema que se propone es el siguiente:

\begin{itemize}
\item Se agrega al mensaje MH\textunderscore DATA los campos: bcast\textunderscore seqnum, para que un nodo pueda diferenciar los mensajes de broadcast; y bcast\textunderscore radio, para establecer cuál es el radio (en términos de reenvío) de alcance del broadcast.

\item Cuando un nodo se prepara para enviar un mensaje de broadcast, debe especificar que este mensaje es de broadcast, asignar un valor de  bcast\textunderscore seqnum coherente y un radio de alcance.

\item El nodo que recibe un mensaje de broadcast debe utilizar el bcast\textunderscore seqnum para decidir si es necesario reenviar o no el mensaje.

\item Una vez que ha decidido reenviar el mensaje, debe decrementar el valor del bcast\textunderscore radio para indicar a los próximos saltos que el mensaje ya ha avanzado un nodo en el radio de alcance.
\end{itemize}

\subsubsection{Encolamiento y reenvío de mensajes}
Para reducir la pérdida de mensajes, se diseña una cola de mensajes que permita encolar los mensajes generados por el nodo y los mensajes recibidos desde otros nodos que deben ser reenviados. Una cola FIFO permite un ordenamiento intrínseco de los mensajes, donde el primer mensaje ingresado estará asociado a la cabeza de la cola FIFO (head) y el último estará en la cola (tail).\\ 	
%Diagrama de encolamiento de mensajes
%\begin{figure}[H]
 %\centering
 %\includegraphics[scale=0.625]{imgs/DiagramaColaFifo.eps}
% \caption{Diagrama de encolamiento de mensajes}
%\end{figure}

Las secuencias, funcionalidades y componentes descritos se implementarán en forma de una Configuración NesC, de forma similar a como se realiza con PEGASIS en el capítulo 3. En el próximo capítulo se describirá como se implementa este diseño para TinyOS, utilizando LPL para realizar las operaciones en ultra bajo consumo. Además se revisan las pruebas de verificación para cada componente y funcionalidad, junto con la evaluación energética correspondiente.
	
	%Incluir capítulo 5
	%	\chapter{Implementación del protocolo con topología de árbol}

\section{Descripción de la implementación}
La implementación del esquema de recolección propuesto en la sección anterior, se realiza programando un ''Componente de recolección'', se divide en dos módulos:  el ''Componente de ruteo'', que es el encargado de formar la red y establecer la ruta hacia el punto de recolección; y la implementación del componente de reenvío, que es el encargado de encolar y reenviar los mensajes.\\

\subsection{Componente de Recolección - CollectionC}
La implementación del componente de recolección, diseñado en capítulo 3, comienza con la programación de una configuración nesC que provee una interfaz similar a la que provee CTP\cite{TEP123}. La configuración se denomina CollectionC, y provee la implementación de las interfaces: \textbf{SplitControl, Send, Receive y Broadcast}. La descripción  de CollectionC se realiza en el archivo CollectionC.nc.

\begin{figure}[H]
	\centering
 	\includegraphics[scale=0.5]{imgs/DiagramaComponenteCollectionC.eps} 
 	\caption{Diagrama de componentes de la configuración CollectionC}
\end{figure}

En la figura 5.1 se describe la configuración utilizando el lenguaje visual de descripción de componentes y sus relaciones, propuesto en el libro TinyOS Programming\cite{TinyOSProgramming}. La configuración consta de dos grandes módulos: El \textbf{módulo RouterP y el módulo ForwardeP}. El módulo RouterP se encarga de incorporar el nodo a la red, proveer al módulo ForwarderP del próximo salto, y tolerar fallas debido a la desaparición de nodos de la red. El módulo ForwarderP se encarga de: colocar en la cola de envío los mensajes que la aplicación necesita enviar (interfaz Send); recibir los mensajes de los nodos hijos; encolar los mensajes recibidos; y detectar la desaparición de nodos. \\

Para que el usuario pueda configurar los parámetros de funcionamiento del protocolo se provee el archivo Config.h, donde, mediante definiciones para el preprocesador C, el usuario puede configurar parámetros que se describirán en esta sección.

\subsection{Componente RouterP}
%Qué es?
El componente Router es el que implementa las secuencias de asociación y reconexión. El componente provee a la configuración la interfaz SplitControl y al módulo ForwarderP la interfaz RouterInfo.

%Interfaz SplitControl-Secuencia de Asociación Beacon Request
\subsubsection{Interfaz SplitControl}
La implementación comienza con la programación de la interfaz SplitControl. Esta interfaz permite a la capa de aplicación, iniciar el componente de ruteo utilizando el comando SplitControl.start(). Cuando se invoca este comando, este inicia el componente LPL con el intervalo definido como LPL\textunderscore INTERVAL en el archivo Config.h. A continuación, inicia la pila ActiveMessage y el timer BeaconRequestTimer con el intervalo BEACON\textunderscore REQUEST\textunderscore INTERVAL, configurable en Config.h. De esta forma se inicia la secuencia de asociación y durante este periodo se mantiene el estado del nodo en la red como NO\textunderscore NWK (sin red). \\

Durante el tiempo que transcurre entre cada BEACON\textunderscore REQUEST, el nodo almacena los BEACON\textunderscore RESPONSE, y cuando se gatilla nuevamente el timer, determina el mejor candidato para realizar la asociación. Si el nodo ha recibido mensajes BEACON\textunderscore RESPONSE y ha determinado un candidato para solicitar la asociación a la red, envía un JOIN\textunderscore REQUEST al candidato. Si no ha encontrado ningún candidato, vuelve a enviar un BEACON\textunderscore REQUEST.\\

Se implementan dos formas de determinar el candidato para requerir asociación: Por mayor RSSI, elige el nodo ''más cercano electromagnéticamente''; y por menor profundidad, se elige el nodo de menor profundidad en la red, pero que tenga un RSSI mayor a -87[dBm], que se ha medido como el mínimo para garantizar una tasa de recepción de mensajes correctos de 85\%  \cite{RSSIUnder}. La utilización de uno de estos esquema se deja configurable en el archivo Config.h, mediante las definiciones SELECT\_BEACON\_BY\_RSSI y SELECT\_BEACON\_BY\_DEPTH. \\

Cuando el nodo recibe un mensaje JOIN\textunderscore OK el nodo establece sus variables de red depth (profundidad) y parent (padre), con los valores obtenidos de este mensaje. Posteriormente cambia al estado IN\_NWK, activando el envío de ACKs con la interfaz CC2420Config y señalizando el evento SplitControl.startDone(SUCCESS) a la aplicación. La señalización de este evento permite que la aplicación esté en conocimiento del momento en que dispone de red para poder enviar mensajes de datos al Gateway.\\

Un punto importante en la implementación, es la incorporación de un sistema de ''Despliegue forzado de la red'', similar al desarrollado para la implementación de PEGASIS, para permitir el desarrollo de pruebas multihop en el laboratorio. El sistema utiliza la matriz ''hopsMatrix'', que implementada en un arreglo bidimensional que describe con un 0, 1 ó 10 la posibilidad de comunicación entre dos nodos. 

\begin{figure}[H]
	\centering
 	\includegraphics[scale=0.5]{imgs/hopsmatrix.eps} 
 	\caption{Ejemplo de matriz utilizada para el sistema de despliegue forzado }
\end{figure}

En la figura 5.2. se muestra una de las matrices utilizada para las pruebas. Cero se utiliza sólo para representar la relación entre un nodo con sí mismo; un 1 en una fila M y columna N, indica que el nodo M puede recibir mensajes del nodo N; mientras que un 10 en la fila M y columna N, indica que el nodo M no puede recibir mensajes del nodo N. Por ejemplo: el 1 en la posición (0,1) indica que el nodo 0 puede recibir mensajes del nodo 1; mientras que el 10 en la posición (0,2), indica que el nodo 0 no puede recibir mensajes del nodo 2. Cuando está activado el despliegue forzado, se utiliza la función ''bool isVisible(uint16\_t id)'' para revisar si un mensaje recibido se procesará o no. Para activar el despliegue forzado, basta con definir ''FORCED\_DEPLOYMENT'' en el archivo Config.h.

\subsection{Componente ForwarderP}
El componente ForwarderP es el encargado de implementar las secuencias de envío y reenvío multihop. Provee a la configuración las interfaces Init, Send, Receive y ForwarderInfo. ForwarderInfo es la interfaz desarrollada para proveer información desde el Forwarder a otros componentes.
%Interfaz Init?

\subsubsection{Interfaz Send}
La implementación del componente Forwarder comienza con la implementación de la interfaz Send. La interfaz Send proporciona el comando Send.getPayload, el comando Send.send, y el evento Send.sendDone.\\

El comando Send.getPayload se utiliza para proporcionar a la aplicación el espacio de memoria para que escriba el mensaje datos que desea enviar. Este comando recibe como entrada un puntero a una estructura ''message\textunderscore t''  y el tamaño del mensaje. Con estos datos, el comando se encarga de crear el espacio de memoria para el mensaje, del tamaño solicitado por la aplicación más el tamaño de la cabecera  ''MHData'', según lo definido en la sección 4.3.1. Finalmente, retorna un puntero a la dirección del payload del mensaje. Esta implementación, permite a la aplicación la utilización de mensajes de largo variable y no limitar a la aplicación a utilizar mensajes de un largo fijo.\\

El comando Send.send es el que permite a la aplicación colocar un mensaje en la cola de envío. Recibe como entrada un puntero a la estructura del mensaje ''message\textunderscore t'' y el largo. 

\begin{figure}[H]
	\centering
 	\includegraphics[scale=0.6]{imgs/DiagramaFlujoComandoSendSend.eps} 
 	\caption{Diagrama de flujo de la implementación del comando Send.send}
\end{figure}

Como muestra el diagrama de flujo en la figura 5.3, el comando Send.send comienza llenando la cabecera del mensaje, utilizando la interfaz RouterInfo para obtener el padre del nodo y la profundidad a la cuál se encuentra en la red. Además, escribe en la cabecera  el campo originaddr con el TOS\textunderscore NODE \textunderscore ID del nodo, y el contador de mensajes de datos en el campo seqnum. Posteriormente, se utiliza la interfaz MessagePool para obtener un espacio de memoria para el mensaje utilizando el comando MessagePool.get(), y si la ejecución del comando resulta exitosa, se copia el mensaje desde la aplicación a la MessagePool. A continuación, se encola el puntero al mensaje almacenado en la MessagePool en la cola FIFO SendQueue, y si resulta exitoso, se encola la tarea sendTask.\\

\begin{figure}[H]
	\centering
 	\includegraphics[scale=0.6]{imgs/DiagramaFlujoSendTask.eps} 
 	\caption{Diagrama de flujo de la implementación de la tarea SendTask}
\end{figure}

%Hablar de la sendTask
La tarea sendTask, es una tarea nesC que se encarga de sacar los mensajes de la cola y enviarlos por el aire. Cada vez que se encola un mensaje, se realiza un post de la tarea sendTask. Si la cola no está vacía, la tarea extrae la cabeza de la cola (el mensaje encolado más antiguo) e intenta enviarlo por el aire, utilizando la pila ActiveMessage. Si el envío se realiza con éxito, el sistema invoca al manejador del evento sendDone (ver figura 5.5), si se ha recibido un ACK del nodo destinatario, el manejador del evento elimina el mensaje más antiguo en la cola y si la cola aún no está vacía, vuelve a realizar un post de la sendTask. Si no se ha recibido un ACK, se vuelve a realizar un post de la sendTask, pero sin eliminar el mensaje más antiguo de la cola, forzando su reintento.

\begin{figure}[H]
	\centering
 	\includegraphics[scale=0.6]{imgs/DiagramaFlujoMHSendDone.eps} 
 	\caption{Diagrama de flujo de la implementación del evento MHSend.sendDone}
\end{figure}

%Reenvío y enconlado de mensajes recibidos
\subsubsection{Interfaz MHReceive.receive}
El componente ForwarderP es el que implementa la interfaz MHReceive.receive, que es la que maneja el evento gatillado por la recepción de mensajes de datos multihop. 

\begin{figure}[H]
	\centering
	\includegraphics[scale=0.6]{imgs/DiagramaFlujoReceiveBroadcast.eps}
	\caption{Diagrama de flujo de la implementación del evento MHReceive.receive}
\end{figure}

Como se muestra en la figura 5.6, cuando se recibe un mensaje de datos, primero se revisa el campo ''type'' en la cabecera para revisar si es un mensaje broadcast. Si lo es, se revisan los campos ''originaddr'' y ''bcast\_seqnum'' para revisar si es un mensaje repetido, si no lo es, se actualiza la cabecera y se señaliza el evento Broadcast.Receive entregando a la aplicación el mensaje, el payload y el largo. Luego de que se señaliza el Broadcast recibido, el mensaje se vuelve a encolar (sólo si el bcast\_radius es mayor que 0 después de ser decrementado). Si el mensaje no era de tipo broadcast, entonces se encola en la cola FIFO SendQueue y se realiza un post de la tarea sendTask.
	
\subsubsection{Interfaz Broadcast}
La interfaz Broadcast se ha implementado con el objetivo de proveer al sistema un medio de comunicación hacia los nodos, que permita por ejemplo: configurar la frecuencia con la que realizan las mediciones, o el sensor del cuál se desean obtener mediciones. También se permite que un nodo envíe un broadcast a los nodos de la red. \\

Los procesos realizados para enviar un mensaje de broadcast son los mismos que se describen en la figura 5.4, de la implementación del comando Send. Difieren en que al llenar la cabecera, el campo ''type'' indica que es un mensaje de broadcast, el ''bcast\_seqnum'' se llena con un valor válido (se incrementa con cada envío broadcast) y también se llena el campo ''bcast\_radius'', con el radio que ha indicado la aplicación. Los procesos de reenvío de los mensajes de broadcast se describieron en la sección anterior.

\section{Pruebas de verificación}
%Prueba de esta funcionalidad
Las pruebas de verificación de la funcionalidades implementadas, se realizan utilizando el sistema de despliegue forzado y la aplicación WSNSniffer. Esta última se modificó para que soportara el formato de los nuevos mensajes especificados en la sección 4.3.1.
 
\subsection{Pruebas al módulo RouterP}
En esta sección se describen las pruebas realizadas y los resultados para verificar las siguientes funcionalidades del módulo RouterP: Descubrimiento de red, Secuencia de Asociación, Asociación por RSSI, Asociación por Profundidad.

\subsubsection{Descubrimiento de red}
La primera prueba consiste en verificar que el mensaje beacon request se envíe periódicamente mientras el nodo no encuentre un nodo en la red con el cuál asociarse. Para esta prueba se utiliza la aplicación de prueba SampleApp, que se encuentra en el CD adjunto a este trabajo, que periódicamente realiza una medición con el sensor de luz y la envía hacia el Gateway utilizando el protocolo implementado. Para esta prueba el TMoteSky que corre SampleApp tiene un ID igual a 1 y un intervalo BEACON\textunderscore REQUEST\textunderscore INTERVAL igual a 1[s].

\begin{table}[H]
\centering
\caption{Extracto del reporte del sniffer para la prueba de descubrimiento de red}
 \includegraphics[scale=0.5]{imgs/ResultadosProcesadosBeaconLoop.eps}
\end{table}

En el cuadro 5.1 se muestra el resultado de la prueba registrado por el sistema Sniffer. La columna ''Timestamp'' muestra el tiempo en que la aplicación Sniffer recibió el mensaje. En promedio, el tiempo de recepción del mensajes es de 975[ms], lo cuál está un 2,5\% por debajo del valor configurado. La columna ''Src Addr'' muestra el TOS\textunderscore NODE\textunderscore ID del dispositivo origen del mensaje, este corresponde al valor 1, que fue el valor configurado para el dispositivo utilizado para esta prueba. La columna ''Dest Addr'' muestra el TOS\textunderscore NODE\textunderscore ID del destinatario del mensaje, en este caso, corresponde a la dirección de broadcast 65535 o 0xFFFF en notación hexadecimal. Ambos identificadores son obtenidos de la cabecera Active Message, para diferenciarlos de los identificadores que se incluyen en la cabecera de los mensajes implementados para la red. La columna ''AM Type'' muestra el identificador Active Message que se ha definido para este tipo de mensajes, que corresponde a 32 (0x21 en hexadecimal) para los mensajes de Topología.\\

Se puede apreciar que tal como se describe la secuencia de asociación en el capítulo 3, el nodo envía beacon request periódicamente , aproximadamente cada 1[s] y se mantiene así al no existir una red en el radio de alcance. 

\subsubsection{Secuencia de Asociación}
A continuación se realiza la misma prueba, pero agregando el nodo Gateway, que por si solo se encuentra en la red, y que debería responder los beacon request del nodo de prueba.

\begin{table}[H]
\centering
\caption{Resultado de prueba de la secuencia de asociación}
 \includegraphics[scale=0.6]{imgs/ResultadosProcesadosTestsJoinRequest.eps}
\end{table}

%Comentario del resultado de la prueba
En el cuadro 5.2 se muestra el resultado registrado por el sistema sniffer para la prueba de secuencia de asociación satisfactoria. En el cuadro se observa que el dispositivo con TOS\textunderscore NODE\textunderscore ID 1 envía un mensaje con subtype = 0x0, que corresponde al mensaje de topología BEACON\textunderscore REQUEST\textunderscore , a la dirección de broadcast. El dispositivo de TOS\textunderscore NODE\textunderscore ID 0 (nodo Gateway) responde con tres mensajes con subtype igual a 0x1, que corresponde al mensaje de topología BEACON\_RESPONSE. El largo de la ráfaga de mensajes de respuesta es configurable en definición del  BEACON\textunderscore RESPONSE\textunderscore BLAST\textunderscore SIZE en el archivo Config.h, esto permite al nodo 1 obtener varias muestras de RSSI y obtener un promedio de estos valores, que le permitirán decidir cuál nodo elegirá como candidato de padre en la red. En cada uno de los mensajes BEACON\_RESPONSE, el campo destaddr igual a 0x0 indica que el nodo 0 está respondiendo al nodo 1. El campo depth igual a 0 indica que la profundidad del nodo 0 en la red es 0. El campo joining igual a 0x1 indica que el nodo está incorporando nodos hijos a la red. El campo isCollector igual a 0x1 indica que el nodo 0 es el recolector (y Gateway). A continuación, el nodo 1 responde al BEACON con un mensaje de subtype igual a 0x2, que corresponde al mensaje de topología JOIN\textunderscore REQUEST. El campo destaddr igual a 0x0 indica que el mensaje está dirigido al nodo 0, ya que el nodo eligió al nodo 0 como candidato de padre en la red. El campo  joining igual a 0x1, esto indicaría que el nodo 1 está aceptando asociaciones, pero ya que no está en la red, no responderá a mensajes BEACON\textunderscore REQUEST\textunderscore , ni el nodo 0 utilizará esa información en el proceso de asociación. El campo isCollector es igual a 0x0, indicando que el nodo 1 no es recolector. El campo redundancy es un campo que se utilizó de manera experimental y se mantuvo en la estructura del mensaje para futuros desarrollos.\\

\subsubsection{Verificación de asociación por mayor RSSI}
Para verificar la funcionalidad de asociación mediante la utilización del mayor RSSI, se realiza el siguiente experimento. Primero se despliegan espacialmente 3 nodos: nodo Gateway (0), nodo 1 y nodo 2. como se muestra en la figura 5.7

\begin{figure}[H]
	\centering
 	\includegraphics[scale=0.6]{imgs/TopologiaDesplieguePruebaAsociacionRssi.eps} 
 	\caption{Topología forzada para prueba de limitación de hijos}
\end{figure}

Como muestra la figura 5.7, los nodos se despliegan linealmente, donde el nodo 1 se encuentra a 3 [m] del nodo 0, mientras que el nodo 2 a 3.5[m] del nodo 0 y 0.5 del nodo 1. Con esta disposición, el nodo 2 debería medir diferencias entre los indicadores RSSI medidos para los mensajes del nodo 1 y el nodo 0.  El objetivo, es que el nodo 2 reciba beacons de ambos nodos, pero que elija el nodo que envió beacons con un RSSI promedio más alto. 

\begin{center}
\begin{table}[H]
\caption{Resultado de prueba de asociación por mayor RSSI}
 \includegraphics[scale=0.6]{imgs/ResultadosProcesadosTestsAsociacionRSSI.eps}
\end{table}
\end{center}

En el cuadro 5.3 se muestra la secuencia de mensajes capturada por el sniffer para la prueba. La secuencia capturada es la secuencia de asociación del nodo 2. Se observa que el nodo 2 en el tiempo 0:46:56:126 inicia la asociación enviando un BEACON\textunderscore REQUEST a la dirección de broadcast. Los nodos 0 y 1 responden con los respectivos mensajes BEACON\textunderscore RESPONSE. Luego, el nodo 2 elige unirse al nodo 1 y en el tiempo 00:46:57:083 envía un mensaje de JOIN\textunderscore REQUEST. Finalmente, recibe un JOIN\textunderscore OK para confirmar la asociación. El resultado de esta secuencia es la asociación del nodo 2 al nodo 1, que es lo esperado si es que el nodo 2 midió un RSSI promedio más alto de los mensajes BEACON\textunderscore RESPONSE del nodo 1 que del nodo 2. Sin embargo, con el sniffer no se puede confirmar ese comportamiento. Por esta razón, se implementó que el nodo 2 enviara por comunicación serial la tabla ' beaconList ', que es el arreglo bidimensional que utiliza para almacenar los datos de los BEACON\textunderscore RESPONSE recibidos.\\

\begin{center}
\begin{table}[H]
\caption{Output recibido por puerto serial del nodo 2}
 \includegraphics[scale=0.6]{imgs/OutputAsociacionRssi.eps}
\end{table}
\end{center}

En el cuadro 5.4 se muestran las 15 entradas, de las cuales sólo las primeras seis son válidas (campo valid igual a 65534 es válido y 65535 es inválido). Estas entradas válidas corresponden a seis mensajes BEACON\textunderscore RESPONSE que recibió el nodo 2. Las entradas 0, 1 y 2 corresponden a mensajes del nodo 0 y el promedio entero de sus RSSI es -32 [dBm]. Las entradas 3, 4 y 5 corresponden a mensajes del nodo 1 y el promedio entero de sus RSSI es 8 [dBm]. De esta manera, se verifica que el nodo 2 eligió como candidato para la asociación al nodo del que recibió los mensajes BEACON\textunderscore RESPONSE con mayor RSSI.\\

\subsubsection{Verificación de asociación por menor profundidad}
%Probar formas de asociación, DEPTH
Para verificar la funcionalidad de la asociación por menor profundidad, se utiliza el mismo despliegue no forzado utilizado y mostrado en la figura 5.7. Además, se vuelve a utilizar la herramienta ''printf'' para enviar por comunicación serial al computador la tabla 'beaconList'.

\begin{center}
\begin{table}[H]
\caption{Resultado de prueba de asociación por menor profundidad}
 \includegraphics[scale=0.6]{imgs/ResultadosProcesadosTestsAsociacionPorDepth.eps}
\end{table}
\end{center}

La secuencia de asociación comienza en el tiempo 0:38:56:207, cuando el nodo 2 envía un mensaje BEACON\textunderscore REQUEST a la dirección de broadcast. A continuación, los nodos 0 y 1 responden con los mensajes BEACON\textunderscore RESPONSE correspondientes. Luego, en el tiempo 0:38:57:164 el nodo 2 solicita una asociación al nodo 0, y este confirma la asociación en el tiempo 0:38:57:200. La secuencia resultante es la esperada de acuerdo al criterio de asociación por menor profundidad, debido a que el nodo 2 decide solicitar la asociación al nodo 0 en vez del nodo 1, ya que el nodo 0 tiene la menor profundidad. Sin embargo, para asegurar que el nodo no está siguiendo el criterio de RSSI, se revisa la tabla de beacons enviada por comunicación serial por el nodo al computador.

\begin{center}
\begin{table}[H]
\caption{Output recibido por puerto serial del nodo 2}
 \includegraphics[scale=0.6]{imgs/OutputAsociacionProfundidad.eps}
\end{table}
\end{center}

Como se observa en el cuadro 5.6, el nodo 1 tiene RSSI mucho mayores que el nodo 0, por lo que claramente el nodo 2 no utilizó el criterio de mayor RSSI para elegir el nodo al cual solicitar la asociación.

\subsection{Pruebas al módulo ForwarderP}
En esta sección se describen las pruebas realizadas y los resultados para verificar las siguientes funcionalidades del módulo ForwarderP: Verificación de reenvío de mensajes y Verificación de broadcast de red.

\subsubsection{Verificación de reenvío de mensajes}
Para verificar que el reenvío de los mensajes va de padre en padre hasta el Gateway, se revisa el reporte del sniffer para un despliegue forzado de la red. La primera prueba, se realiza con el despliegue descrito en la figura 5.8.

\begin{figure}[H]
	\centering
 	\includegraphics[scale=0.6]{imgs/DesplieguePruebaToleranciaFallas.eps} 
 	\caption{Despliegue para primera prueba de verificación de reenvío}
\end{figure}
 
En el cuadro 5.7 se muestra un Extracto del reporte del sniffer entregado para esta prueba. Para verificar que los mensajes viajan de padre en padre hasta el Gateway (nodo 0), basta con observar como viaja, por ejemplo, el mensaje generado por el nodo 4 en el timestamp 21:55:54:790. Se observa en la columna SrcAddr que ese mensaje fue enviado en ese salto por el nodo 4, y en la columna Data, que la dirección de origen del mensaje de datos es el nodo 2. Por lo tanto, ese mensaje fue generado inicialmente por el nodo 4. A continuación, se observa que el mensaje reenviado por el nodo 3 al 1, es exactamente el mismo que envió anteriormente el nodo 4. Finalmente, en el timestamp 21:55:55:045 el mismo mensaje es reenviado a su destino final, el Gateway (nodo 0). Si se realiza el mismo ejercicio con el resto de los mensajes, se confirma que todos los mensajes generados por los 5 nodos de la red, fueron reenviados de padre en padre hasta el Gateway. 
 
\begin{center}
\begin{table}[H]
\caption{Resultado de la primera prueba de reenvío}
 \includegraphics[scale=0.6]{imgs/ResultadosProcesadosTestReenvioTop2.eps}
\end{table}
\end{center}
%Probar broadcast
\subsubsection{Verificación de broadcast de red}
Para verificar la funcionalidad de la implementación del broadcast de red, se realiza el siguiente experimento:

\begin{itemize}
 \item Se implementa utilizando la configuración de despliegue forzada, una red lineal, con 4 nodos, donde el nodo 1 es el más cercano al Gateway y el 4 es el más lejano.

 \item Se implementa una aplicación de prueba que envía un mensaje de broadcast, con radio dinámico. Cada vez que se presiona el botón de usuario, con esto gatilla el envío del mensaje de broadcast, y el incremento del radio para el próximo mensaje.

 \item Una vez formada la red, se presiona el botón del nodo 4 varias veces y se captura con el Sniffer la secuencia de mensajes generada.
\end{itemize}

En el cuadro 5.8, se observa el resultado capturado por el sniffer cuando el nodo 4 envía el broadcast con radio 2. El mensaje es enviado por el nodo 4 en el timestamp 19:54:45:245, y a continuación es reenviado por el resto de los nodos, y con cada salto, se decrementa el campo ''bcast\_radius''.

\begin{center}
\begin{table}[H]
\caption{Secuencia capturada del broadcast con radio 2}
 \includegraphics[scale=0.6]{imgs/ResultadosProcesadosTestsBroadcastRadio2.eps}
\end{table}
\end{center}
%En el cuadro 5.9, se observa cuando el nodo envía el broadcast con radio 3.
%\begin{center}
%\begin{table}[H]
 %\includegraphics[scale=0.6]{imgs/ResultadosProcesadosTestsBroadcastRadio3.eps}
%\caption{Secuencia capturada del broadcast con radio 3}
%\end{table}
%\end{center}

%El mensaje es enviado por el nodo 4 en el timestamp 19:56:02:723, y a continuación es reenviado por el resto de los nodos, y con cada salto, se decrementa el campo ''bcast\_radius''.
\subsection{Prueba de verificación de tolerancia a fallas}
Para verificar la tolerancia a falla para la desaparición de un nodo se utiliza el despliegue forzado que se muestra en la figura 5.9.

\begin{figure}[H]
	\centering
 	\includegraphics[scale=0.6]{imgs/DesplieguePruebaToleranciaFallas.eps} 
 	\caption{Despliegue la prueba de desaparición de un nodo}
\end{figure}

Con la red funcionando, se remueven las baterías del nodo 3 y se registra el comportamiento con el sniffer.\\

En el cuadro 5.9 se muestra el resultado obtenido. Como se ha configurado en el archivo Config.h, el nodo 4 intenta enviar tres veces el mensaje de datos con seqnum igual a 0x5 al nodo 3. Luego de los reintentos, en el timestamp 22:29:05:902 el nodo envía un BEACON\_REQUEST para solicitar una nueva asociación. El nodo 2, responde satisfactoriamente y asocia al nodo 4 nuevamente. El nodo 5 realiza la misma operación, el registro completo de la prueba se encuentra en el CD adjunto a este trabajo.\\

\begin{table}[H]
\centering
\caption{Registro de la detección del nodo perdido y la reconexión del nodo 4}
 \includegraphics[scale=0.6]{imgs/ResultadosProcesadosTestsToleranciaFallaMuerte.eps}
\end{table}

\subsection{Prueba de verificación de red multihop}
%Prueba Depto elo.
Para validar el funcionamiento multihop en un escenario ''medianamente real'' se desplegó la red en las instalaciones del departamento de electrónica (Edificio B, 3er piso), por un periodo tiempo de 30 minutos. En el archivo Config.h, se deshabilita la opción de despliegue forzado y se utilizó el sistema de asociaciones por RSSI  En la figura 5.10 se muestra la disposición en el plano de los nodos.

\begin{figure}[H]
 \centering
 \includegraphics[scale=0.5]{imgs/RedDeptoElo.eps}
 \caption{Despliegue de los nodos en el Departamento de Electrónica. Tercer piso, edificio B, Casa Central de la UTFSM}
\end{figure}

Los resultados fueron bastante satisfactorios, la formación de la red funcionó sin problemas, estableciendo la siguiente topología:

\begin{table}[H]
\centering
\caption{Asociaciones resultantes en el despliegue en el Departamento de Electrónica}
 \includegraphics[scale=0.6]{imgs/ResultadosProcesadosTestsTopologiaDeptoElo.eps}
\end{table}
 
 Además, analizando el reporte se verifica que se rutearon al Gateway la totalidad de los mensajes generados.
 
\section{Evaluación Energética}
%Utilizando el mismo sistema de medición de corriente consumida presentado en la sección XX, se realiza una evaluación del protocolo en distintos escenarios. Primero, se realiza un experimento para medir el comportamiento energético en función del intervalo LPL y el intervalo de generación de mensajes. A continuación se realiza otro experimento para medir el comportamiento energético en función de la topología de la red.

%\subsection{En función del intervalo de generación de mensajes}

Utilizando el mismo sistema de medición de corriente consumida presentado en la sección 3.4.1, se realiza una evaluación del protocolo, forzando el despliegue para que resulte en una topología de árbol con dos niveles, donde en el primer nivel se encuentra el nodo 1 (hijo del Gateway) y 4 nodos restantes en el segundo nivel (hijos del nodo 1). En los nodos se utilizó la aplicación SampleApp (incluida en el CD adjunto a este trabajo), que periódicamente toma una medición con el sensor de luz y la envía hacia el Gateway, utilizando el protocolo de red implementado. Se utilizó un intervalo LPL fijo de 500 milisegundos, y se midió la corriente media utilizada por el nodo 1 (el nodo que maneja más tráfico) para distintos intervalos de generación de mensajes.

\begin{figure}[H]
 \centering
 \includegraphics[scale=0.4]{imgs/MiProtoTopEstrellaFrecuenciamensajesLpl500.eps}
 \caption{Consumo en nodo con más tráfico en un árbol de dos niveles}
\end{figure}

En la figura 5.11, se observa que la corriente media utilizada por el TMoteSky para los intervalos de 200[ms] y 400[ms] alcanza los niveles más bajos, esto se debe a que esos intervalos son menores al intervalo LPL, y en estos casos la tasa de recepción de mensajes correctos fue nula. Lo anterior sólo se realizó con fines experimentales, para determinar cuál sería el comportamiento para intervalos de envío de mensajes menores al intervalo LPL. Según los resultados obtenidos, el  consumo máximo corresponde a 7,44[mA] con un intervalo de generación de mensajes de 600 [ms]. El consumo mínimo fue para el intervalo de 2000[ms] y equivale a 3,19[mA]. Con este mínimo de consumo, se estima que utilizando baterías que proporciones una carga de 2400[mAh], la autonomía del nodo alcanza los 30 días, cumpliendo con el mínimo de lo propuesto en los objetivos (ver Sección 1.3).\\

Como se muestra en la figura 5.11, posteriormente al intervalo de 600 [ms], el consumo decrece progresivamente. En consecuencia, se realiza la misma evaluación para intervalos de generación de mensajes más largos, como se muestra en la figura 5.12.

\begin{figure}[H]
 \centering
 \includegraphics[scale=0.3]{imgs/GraficoConsumoLpl500IntervaloLargo.eps}
 \caption{Consumo en nodo con más tráfico en un árbol de dos niveles, para intervalos de generación de mensajes más largos}
\end{figure}

Como muestra la figura 12, con intervalos más largos el consumo es aún menor. Con 30000 [ms] se obtiene una corriente media de 0,95 [mA]. Con este consumo, se estima que utilizando baterías que proporcionen una carga de 2400[mAh], la autonomía del nodo alcanza los 105 días, lo que representa más de 3 meses.\\

Para contrastar los resultados con otro protocolo, se implementó una aplicación que utiliza el protocolo CTP para recolectar mediciones. Sin embargo, cuando se activa el componente LPL en esta aplicación, el protocolo no logra formar una red. Este comportamiento ha sido documentado por  el grupo de trabajo de red para TinyOS2.x\cite{Net2WG} y en su programa de trabajo se incluye agregar la compatibilidad de CTP con LPL. En febrero de 2010 se revisó por última vez el código fuente de CTP y aún no se ha incorporado la compatibilidad, según el sitio del grupo de trabajo, todavía está en periodo de pruebas.

%\subsubsection{En función de la topologia}
%Se implementa una red con un árbol de dos niveles, igual a la implementada en la sección XX, pero los nodos del segundo nivel se ingresan de a uno, para medir el consumo de corriente el nodo padre en función del número de hijos.
 
 
 
%Para medir la distribución del consumo de corriente en función de la profundidad, se despliega una red con topología lineal, es decir, un árbol de cinco niveles, con un nodo por nivel.


	
	%Incluir capítulo 5
	%	\chapter{Trabajo futuro y conclusiones}

Luego de completar el trabajo, se tiene una clara visión de las mejoras para una próxima iteración de este trabajo y las conclusiones que se obtienen según los objetivos iniciales planteados.
\section{Trabajo Futuro}

En el protocolo de árbol de recolección implementado en capítulo 5, es posible que se generen círculos en las asociaciones. Estos círculos causarían que los mensajes generados por un nodo en vez de llegar al Gateway, transiten por este círculo y nunca salgan de él. Para una futura revisión del protocolo, sería conveniente implementar un mecanismo para la detección de estos círculos, para dar mayor robustez a la implementación.\\

Si bien la implementación del Protocolo PEGASIS no logró cumplir con los requerimientos de este trabajo, es posible, que aplicar un esquema similar sobre el protocolo de árbol de recolección que se implementó en el capítulo 5, permita un mejor desempeño energético de la red. Esto se debe a que las ideas de utilizar la redundancia (en términos de que más de un nodo puede rutear los mensajes de un nodo hijo, no sólo el padre) para repartir el tráfico en los nodos y la variación dinámica en la potencia de transmisión, podrían optimizar el consumo de corriente de los dispositivos y repartirlo por la red.\\

La aplicación WSNSniffer, que se desarrolló para verificar el comportamiento de la red, fue de gran utilidad. Para una próxima revisión, este programa puede ser llevado a otro nivel, incorporando la capacidad de detectar los mensajes de protocolos como CTP, Dissemination, FTSP y las aplicaciones que se distribuyen como ejemplo en el sistema TinyOS. También, se podría agregar una funcionalidad que permita agregar dinámicamente los mensajes que se desea detectar. Esta herramienta podría ser de gran ayuda para el estudio de las redes de sensores inalámbricos y para futuros desarrollos.

%¿Qué escribio alviña en sus conclusiones?

%Primer párrafo: en el primer parrafo concluye las debilidades de los actuales sistemas de monitoreo.

%Segundo párrafo: Habla sobre el porque de su solucion y que conocimientos aplicó para llevarla a cabo.

%Tercer párrafo: Explica que que soluciona su sistema y qué es lo que hace.

%Cuarto Habla de demora y dificultades en el desarrollo.

%Quinto y último párrafo: Habla de las proyecciones a futuro de su trabajo.
\section{Conclusiones}
%Concluciones acerca de la investigación.
Los algoritmos como LEACH, APTEEN y PEGASIS, no pueden por si solos reducir el consumo energético total de un dispositivo como el TMoteSky, para alcanzar un tiempo de operación de nodos del orden de meses. Para lograr esto, es necesario utilizar otras técnicas como Low Power Listening.\\

Los protocolos de ruteo implementados para redes de sensores inalámbricos que utilizan el TmoteSky y TinyOS (CTP y Dissemination), no están diseñados para cuidar el recurso energético del dispositivo. Actualmente, no están diseñados para funcionar en conjunto con el componente Low Power Listening que es el que permite reducir el consumo de corriente por el problema de ''Idle Listening'', y así, extender el tiempo de operación de los dispositivos al orden de meses. \\

Como se verificó en las secciones 5.2 y 5.3, se implementó satisfactoriamente un protocolo que rutea los mensajes hacia el Gateway a través de un árbol de recolección. El protocolo tolera la desaparición de nodos y utiliza el componente Low Power Listening. Este último permite: desplegar una red de sensores, autónomos energéticamente, en un campo sin distribución eléctrica; los dispositivos permitirían recolectar mediciones, por ejemplo, de humedad y temperatura, realizando esta operación por al menos un mes y hasta 3 meses (de acuerdo con los resultados obtenidos en los escenarios utilizados). La solución propuesta permitiría desarrollar una solución concreta al problema de recolección de datos basado en una red de sensores inalámbricos, utilizando un esquema multi hop, sobre una red asincrónica y permitiendo tiempo de operación de los nodos del orden de 1 a 3 meses, según el hardware energético que se utilice.
	
	%	
\begin{thebibliography}{99}
\begin{sloppypar}

% Javier REFS
\bibitem{alexa-twitter}Alexa - Twitter.com site info. [en línea]\textless \url{http://www.alexa.com/siteinfo/twitter.com}\textgreater[consulta: 30 Octubre 2012] 


%Capitulo 1
\bibitem{WileyWSN} SOHRABY, Kazem. MINOLI, Daniel. ZNATI, Taiebi. Wireless sensor networks: technology, protocols, and applications. Editorial John Wiley \& Sons, Inc. 2007.

\bibitem{WSNSurvey} AL-KARAKI, Jamal N. KAMAL, Ahmed E. Routing Techniques in Wireless Sensor Networks: A Survey. The Hashemite University, Iowa State University. 2004.

\bibitem{ProyectoForestal} Redes inalámbricas de sensores para la detección de incendios forestales y monitoreo de variables de estado de combustibles para el recurso forestal de la Región de Valparaíso. [en línea]\textless \url{http://www.gisincendiosforestales.cl}\textgreater[consulta: 15 junio 2009] 

\bibitem{Wiseconn} Wiseconn - Wireless Sensors. [en línea]\textless \url{http://www.wiseconn.cl/}\textgreater[consulta: 15 junio 2009] 

\bibitem{LeachOriginal} HEINZELMAN, Wendi. CHANDRAKASAN, Anantha. BALAKRISHNAN, Hari. Energy-Efficient Communication Protocol for Wireless  Microsensor Network. Massachusetts Institute of Technology, Cambridge. 2000. [en línea]\textless \url{http://ieeexplore.ieee.org/xpl/freeabs_all.jsp?arnumber=926982}\textgreater[consulta: 20 junio 2008] 

\bibitem{PegasysOriginal} LINDSEY, Stephanie. RAGHAVENDRA, Cauligi. PEGASIS: Power-Efficient Gathering in Sensor Information Systems. Computer Systems Research Department. The Aerospace Corporation. 2003. [en línea]\textless \url{http://ieeexplore.ieee.org/xpl/freeabs_all.jsp?arnumber=1035242}\textgreater[consulta: 20 junio 2008] 

\bibitem{HeedOriginal} YOUNIS, Ossama. FAHMY, Sonia. HEED: A Hybrid, Energy-Efficient, Distributed Clustering Approach for Ad-hoc Sensor Networks. Department of Computer Sciences, Purdue University. 2004. [en línea]\textless \url{http://ieeexplore.ieee.org/xpl/freeabs_all.jsp?arnumber=1035242}\textgreater[consulta:  2 julio 2009] 

\bibitem{TeenOriginal} MANJESHWAR, Arati. AGRAWAL, Dharma. TEEN: A Routing Protocol for Enhanced Efficiency in Wireless Sensor Networks. Center of Distributed and Mobile Computing, ECECS Department, University of Cincinnati, Cincinnati. 2002. [en línea]\textless \url{http://ieeexplore.ieee.org/xpl/freeabs_all.jsp?arnumber=925197}\textgreater[consulta: 20 junio 2008] 

\bibitem{TEP118} LEVIS, Philip. TOLLE, Gilman. TEP 118: Dissemination of Small Values. TinyOS Net2 Working Group. 2007. 
[en línea]\textless \url{http://www.tinyos.net/tinyos-2.x/doc/html/tep118.html}\textgreater[consulta: 21 junio 2008] 

\bibitem{TEP123} FONSECA, Rodrigo. GNAWALI, Omprakash. JAMIESON, Kyle. KIM, Sukun. LEVIS, Philip. WOO, Alec. TEP 123: The Collection Tree Protocol. TinyOS Network Working Group. 2007. [en línea]\textless\url{http://www.tinyos.net/tinyos-2.x/doc/html/tep123.html}\textgreater[consulta: 21 junio 2008] 

\bibitem{IntelMote2Web} ''Intel Mote 2 - Platformx Wiki". [en línea]\textless\url{http://embedded.seattle.intel-research.net/wiki/index.php?title=Intel_Mote_2}\textgreater[consulta: 3 marzo 2010] 

\bibitem{IrisDatasheet} '' Iris Datasheet ''. Crossbow Technology. [en línea]\textless\url{http://www.xbow.com/Products/Product_pdf_files/Wireless_pdf/IRIS_Datasheet.pdf}\textgreater[consulta: 3 marzo 2010] 

\bibitem{Mica2Datasheet} ''Mica2 Datasheet ''. Crossbow Technology.
[en línea]\textless\url{http://www.xbow.com/Products/Product_pdf_files/Wireless_pdf/MICA2_Datasheet.pdf}\textgreater[consulta: 3 marzo 2010] 

\bibitem{MicazDatasheet} ''Micaz Datasheet ''. Crossbow Technology.
[en línea]\textless\url{http://www.xbow.com/Products/Product_pdf_files/Wireless_pdf/MICAz_Datasheet.pdf}\textgreater[consulta: 3 marzo 2010] 

\bibitem{ShimmerWiki} ''Shimmer - TinyOS Documentation Wiki''.
[en línea]\textless\url{http://docs.tinyos.net/index.php/SHIMMER}\textgreater[consulta: 3 marzo 2008] 

\bibitem{TelosBDatasheet} ''TelosB Datasheet ''. Crossbow Technology.
[en línea]\textless\url{http://www.xbow.com/Products/Product_pdf_files/Wireless_pdf/TelosB_Datasheet.pdf}\textgreater[consulta: 3 marzo 2020] 

\bibitem{TinyNodeWeb} ''TinyNode''.
[en línea]\textless\url{http://www.tinynode.com}\textgreater[consulta: 3 marzo 2010] 

\bibitem{ATAVRRZ201} ''Atmel Products - MCU Wireless - ATAVRRZ201''. Atmel Corporation.
[en línea]\textless\url{http://www.atmel.com/dyn/products/tools_card_v2.asp?tool_id=4186}\textgreater[consulta: 3 marzo 2010] 

%Capítulo 2
\bibitem{TelosBConsumo} POLASTRE, Joseph. SZEWCZYK, Robert. CULLER, David. Telos: Enabling Ultra-Low Power Wireless Research. Computer Science Department, University of California, Berkeley. 2005. 
[en línea]\textless\url{http://ieeexplore.ieee.org/xpl/freeabs_all.jsp?arnumber=1440950}\textgreater[consulta: 30 Agosto 2009] 

\bibitem{EmbeddedArchitecture} LEVIS, Philip. Embedded Sensor Architecture. Stanford University. Stanford. 2005.
[en línea]\textless\url{http://csl.stanford.edu/ ˜pal/talks/ee282-f05-sensors.pdf}\textgreater[consulta: 31 Agosto 2009]

\bibitem{LeachOptimo} YANG, Haiming. SIKDAR, Biplab. Optimal Cluster Head Selection in the LEACH Architecture. Department of ECSE, Rensselaer Polytechnic Institute. 2007.
[en línea]\textless \url{http://ieeexplore.ieee.org/xpl/freeabs_all.jsp?arnumber=4197919}\textgreater[consulta: 21 junio 2008] 

\bibitem{mumhr} HAMMOUDEH, Mohammad. KURZ, Alexander. GAURA, Elena. MuMHr: Multi-path, Multi-hop Hierarchical Routing. Cogent Computing Applied Research Centre, Coventry University. 2007.
[en línea]\textless \url{http://ieeexplore.ieee.org/xpl/freeabs_all.jsp?arnumber=4394911}\textgreater[consulta: 14 Octubre 2008] 

\bibitem{mleach} MHATRE, V. ROSENBERG, C. Homogeneous vs Heterogeneus Clustered Sensor Networks: A Comparative Study. 2004.
[en línea]\textless \url{http://ieeexplore.ieee.org/xpl/freeabs_all.jsp?arnumber=1313223}\textgreater[consulta: 14 Octubre 2008]

\bibitem{hpegasis} LINDSEY, Stephanie. RAGHAVENDRA, Cauligi. Data Gathering in Sensor Networks using the Energy*Delay Metric. Computer Systems Research Department, The Aerospace Corporation. 2002.
[en línea]\textless \url{http://ieeexplore.ieee.org/xpl/freeabs_all.jsp?arnumber=925196}\textgreater[consulta: 14 Octubre 2008] 

\bibitem{What2ndGenHolds} LEVIS, Philip. What the Second Generation Holds. Presentación para la Stanford University. 2007.
[en línea]\textless \url{http://www.youtube.com/watch?v=j6hRsue5b30}\textgreater[consulta: 8 Octubre 2008] 

\bibitem{TricklePaper} LEVIS, Philip. PATEL, Neil. CULLER, David. SHENKER, Scott. Trickle: A Self-Regulating Algorithm for Code Propagation and Maintenance in Wireless Sensor Networks. 2004. University of California, Berkeley.

[en línea]\textless \url{http://csl.stanford.edu/~pal/pubs/trickle-nsdi04.pdf}\textgreater[consulta: 16 Junio 2009] 

\bibitem{PaperComparacion} SEON, Choong. THUA, Trong. Delay Efficient Multi-Hop Routing Secheme for 
Wireless Sensor Networks. 2006.
[en línea]\textless \url{http://portal.acm.org/citation.cfm?id=1185006}\textgreater[consulta: 26 Marzo 2009] 

%Capítulo 4
\bibitem{TinyOSProgramming} LEVIS, Philip. GAY, David. TinyOS Programming. Editorial Cambridge University Press. 2009.

\bibitem{FtspBug}PacketTimeStamp CC2420 bug.
[en línea]\textless \url{http://docs.tinyos.net/index.php/PacketTimeStamp_CC2420_bug}\textgreater[consulta: 31 Agosto 2009] 

\bibitem{IEEE802154} IEEE 802.15 WPAN TASK GROUP 4. IEEE Standard for information technology-Telecommunications and information exchange between systems-Local and metropolitan area networks-Specific requirements, Part 15.4: Wireless Medium Access Control (MAC) and Physical Layer (PHY) Specifications for Low-Rate Wireless Personal Area Networks (WPANs). IEEE Computer Society. 2006.

\bibitem{RSSIUnder} SRINIVASAN, Kannan. LEVIS, Philip. RSSI is Under Appreciated. Departmen of Electrical Engineering, Department of Computer Science, Standford University. 2006.
[en línea]\textless \url{http://csl.stanford.edu/~pal/pubs/emnets06.pdf}\textgreater[consulta: 21 noviembre 2009] 

%Capítulo 5
\bibitem{Net2WG} Net2WG - TinyOS 2.x Working Group.
[en línea]\textless \url{http://tinyos.stanford.edu:8000/Net2WG}\textgreater[consulta: 2 noviembre 2009] 

\end{sloppypar}
\end{thebibliography}


	%Cambiar nombre de Apéndice a Anexos
	\renewcommand{\appendixname}{Anexo}

	%Comienzo de anexos
	\appendix
	%Incluir anexo A
    %	
\chapter{Programación en TinyOS}
TinyOS es un sistema operativo para sistemas embebidos de código abierto, diseñado especialmente para desarrollar aplicaciones para redes de sensores inalámbricos. Su desarrollo comenzó en la Universidad de Berkeley a fines de los años 90 y actualmente su desarrollo está a cargo de la TinyOS Alliance. TinyOS está diseñado para permitir concurrencia, operación con recursos limitados, soporte para varias plataformas de hardware, y adaptación a la evolución del hardware.\\

El sistema y sus aplicaciones se desarrollan utilizando el lenguaje de programación nesC. Este lenguaje es un dialecto C, que incorpora conceptos como: componentes, interfaces, comandos, tareas y eventos. El ambiente de desarrollo para este sistema está construido para sistemas Unix, sin embargo, puede funcionar sobre sistemas Windows utilizando la herramienta Cygwin. Dentro de las herramientas que provee TinyOS, se pueden encontrar abstracciones para timers, protocolos de comunicación serial, comunicación por radio, protocolos de ruteo, manejo de sensores y persistencia de datos. Actualmente, el sistema soporta las plataformas de hardware: Intelmote2, Iris, Mica, Mica2, Micaz, Shimmer, Telos, TelosB (equivalente a TmoteSky) y Tinynode. 

\section{Instalación y compilación}
El primer paso para desarrollar aplicaciones para el sistema TinyOS es configurar el ambiente de desarrllo. A continuación, se describirán los pasos para instalar el ambiente en un sistema Windows, y posteriormente, los pasos para un sistema Linux.

\subsection{Sistemas Windows}
Para utilizar TinyOS en un ambiente Windows, es necesario instalar la aplicación Cygwin. Esta herramienta habilita un ambiente Linux sobre Windows para ejecutar aplicaciones desarrolladas para el estándar POSIX. Además, es necesario instalar los compiladores para las distintas plataformas que soporta TinyOS, herramientas específicas de TinyOS y el código fuente del sistema. Finalmente, es necesario configurar las variables de entorno para la compilación del sistema. Este tutorial está basado en el manual oficial de instalación publicado en \url{http://docs.tinyos.net/index.php/Installing_TinyOs_2.1} y ha sido probado utilizando Windows XP SP3. 

\subsubsection{Instalación de JDK 1.6}
Realice la instalación del JDK 1.6, descargando la distribución que más le acomode desde \url{java.sun.com}.

\subsubsection{Instalación de Cygwin}
Para la instalación de la aplicación Cygwin debe seguir los siguientes pasos.

\begin{enumerate}
\item Descargar el archivo cygwin-files.zip desde \url{http://cone.informatik.uni-freiburg.de/people/aslam/cygwin-files.zip} y extraer los archivos en una carpeta local (Por ejemplo, C:/cygwin-files/).
\item Ir a la carpeta local y ejecutar el archivo setup.exe.
\item En la pantalla ''Choose installation type'' elija ''Install from a Local Directory''.

\begin{center}
\begin{figure}[H]
\centering
\includegraphics[scale=1]{imgs/anexo_tinyos_1.eps} 
\caption{Selección de tipo de instalación para Cygwin.}
\end{figure}
\end{center}

\item En la pantalla ''Select Root Install Directory'' utilice la configuración como se muestra en la figura A.2.

\begin{center}
\begin{figure}[H]
\centering
\includegraphics[scale=1]{imgs/anexo_tinyos_2.eps} 
\caption{Selección del directorio raíz para Cygwin.}
\end{figure}
\end{center}

\item En la pantalla "Select Local Package Directoy" seleccione la carpeta donde extrajo los archivos.
\item Complete el resto de la instalación sin modificaciones.
\end{enumerate}

\subsubsection{Instalación de compiladores}
Para realizar la instalación de las herramientas de compilación, debe seguir los siguientes pasos.

\begin{enumerate}
\item Una vez instalado Cygwin, ejecute el acceso directo creado en su escritorio o en el menú de inicio para abrir un terminal de Cygwin.

\begin{center}
\begin{figure}[H]
\centering
\includegraphics[scale=0.80]{imgs/anexo_tinyos_3.eps} 
\caption{Terminal Cygwin.}
\end{figure}
\end{center}

\item En esta pantalla usted podrá ejecutar varios de los comandos de uso habitual del ambiente Unix, en este terminal podrá realizar la instalación de los paquetes RPM que contienen los instaladores para las plataformas soportadas por TinyOS.

\item Descargue los paquetes RPM para el MSP430 desde \url{http://docs.tinyos.net/index.php/Installing_TinyOS_2.1.1#Manual_installation_on_your_host_OS_with_RPMs}. Debe descargar todos los paquetes para Windows que se encuentran en la sección Step 3: Install native compilers, TI MSP430 Tools.

\item A continuación, copie los paquetes descargados a una carpeta dentro de la carpeta ''C:/cygwin/home/NombreDeUsuario'' (donde ''NombreDeUsuario'' es el nombre del usuario con el cuál realizó la instalación de Cygwin).

\item Para instalar cada paquete, debe ejecutar el comando ''rpm -ivh nombre\_del\_paquete.rpm''. En caso de error siga las siguientes instrucciones:

\begin{itemize}
\item Si recibe el mensaje ''warning: user TinyOS User does note exist -using root'', ignórelo. 
\item Si recibe el mensaje ''package nombre is intended for a cygwin\_nt-5.1 ... '', vuelva a intentar la instalación ejecutando ''rpm -ivh --ignoreos nombre.rpm''.
\item Si recibe el mensaje ''error: Failed dependencies...'', vuelva a intentar la instalación ejecutando ''rpm -Uvh --force --nodeps nombre.rpm''.
\item Si recibe un mensaje acerca de conflictos con paquetes instalados previamente, vuelva a intentar la instalación ejecutando ''rpm -Uvh --ignoreos --force nombre.rpm''.
\end{itemize}
\end{enumerate}

\subsubsection{Instalación de herramientas específicas}
Para instalar las herramientas específicas siga los siguientes pasos.

\begin{enumerate}
\item Descargue todos los paquetes para Windows desde  \url{http://docs.tinyos.net/index.php/Installing_TinyOS_2.1.1#Manual_installation_on_your_host_OS_with_RPMs}, de la sección Step 4: Install TinyOS toolchain, TinyOS-specific Tools.

\item Instale el paquete nesc-1.3.0-1.cygwin.i386.rpm utilizando el comando ''rpm -Uvh'', si falla, intente con ''rpm -Uvh --ignoreos''.
\item Instale el paquete tinyos-deputy-1.1-1.cygwin.i386.rpm utilizando el comando ''rpm -i'', si falla, intente con ''rpm -i --ignoreos''.
\item Instale el paquete tinyos-tools-1.3.0-1.cygwin.i386.rpm utilizando el comando ''rpm -ivh''.
\end{enumerate}

\subsubsection{Instalación del código fuente de TinyOS y configuración del entorno}
Descargue el paquete tinyos-2.1.0-2.cygwin.noarch.rpm desde  \url{http://www.tinyos.net/dist-2.1.0/tinyos/windows/tinyos-2.1.0-2.cygwin.noarch.rpm}. A continuación, debe instalarlo ejecutando el comando ''rpm -ivh nombre\_del\_paquete.rpm''. A continuación, cree un archivo llamado tinyos.sh y editelo colocando el siguiente código:

\lstinputlisting[language=bash]{tinyos.sh}

Para finalizar, coloque tinyos.sh en la carpeta ''C:/cygwin/etc/profile.d''.

\subsubsection{Instalación del Driver USB Serial FTDI}
Para que el sistema pueda comunicarse con el Tmote Sky, es necesario instalar el driver FTDI que crea un puerto COM virtual para comunicar el dispositivo mediante el puerto USB. Puede encontrar el driver en el sitio de FTDI \url{http://www.ftdichip.com/Drivers/VCP.htm}.

\subsubsection{Verificar instalación del entorno de desarrollo}
Abra un terminal Cygwin y vaya a la carpeta de la aplicación de ejemplo ''Blink'' , ejecutando ''cd /opt/tinyos-2.x/apps/Blink''. En esta carpeta, compile la aplicación para el dispositivo Tmote Sky ejecutando "make tmote", a continuación se presenta la salida de una compilación correcta.

\begin{center}
\begin{figure}[H]
\centering
\includegraphics[scale=0.60]{imgs/anexo_tinyos_4.eps} 
\caption{Resultado de la compilación exitosa de la aplicación Blink.}
\end{figure}
\end{center}

Para verificar que Cygwin se puede comunicar con el dispositivo, conecte el Tmote y ejecute "motelist". Si el driver USB serial está correctamente instalado debe ver una salida similar a la que se muestra a continuación.

\begin{center}
\begin{figure}[H]
\centering
\includegraphics[scale=0.60]{imgs/anexo_tinyos_5.eps} 
\caption{Resultado exitoso de la ejecución del comando motelist.}
\end{figure}
\end{center}

Ahora, para verificar que puede cargar programas correctamente en el  dispositivo, ejecute make tmote install (en el mismo directorio y con el dispositivo conectado). Si la carga del programa se realizó correctamente debería ver una salida similar a la que se muestra en la figura A.6.

\begin{center}
\begin{figure}[H]
\centering
\includegraphics[scale=0.60]{imgs/anexo_tinyos_6.eps} 
\caption{Resultado exitoso de la carga de una aplicación en el TMoteSky.}
\end{figure}
\end{center}

Luego de realizar este tutorial, su sistema Windows debería quedar configurado para desarrollar aplicaciones para TinyOS 2.1.

\subsection{Sistemas Linux}
Para utilizar TinyOS en un sistema Linux (este tutorial está dirigido a la distribución Ubuntu) se puede realizar una instalación automática de las herramientas, utilizando los repositorios y el sistema de paquetes Debian. 

\subsubsection{Agregar el repositorio de TinyOS para Ubuntu}
Para agregar el repositorio a Ubuntu, es necesario abrir con un editor de textos el archivo etc/apt/sources.list. A continuación, hay que agregar la siguiente línea al final del archivo.

\begin{verbatim}
 deb http://tinyos.stanford.edu/tinyos/dists/ubuntu <distribution> main
\end{verbatim}

Las distribuciones soportadas hasta el momento son: edgy, feisty, gutsy, hardy, jaunty, karmic, lucid. Para finalizar este paso, es necesario actualizar el caché utilizando el comando que se muestra a continuación.

\begin{verbatim}
 sudo apt-get update
\end{verbatim}

A continuación, para instalar tinyos y sus herramientas se debe ejecutar el siguiente código.

\begin{verbatim}
 sudo apt-get install tinyos
\end{verbatim}

Es posible que en el repositorio existan varias versiones de TinyOS, si el comando anterior falla y le sugiere versiones, especifique la versión que desea instalar. Para ver las versiones, puede presionar tab luego de escribir tinyos en el comando anterior. Posteriormente, debe agregar el siguiente código al final del archivo .bashrc.

\begin{verbatim}
if [ -f /opt/tinyos-2.1.0/tinyos.sh ] ; then
. /opt/tinyos-2.1.0/tinyos.sh
fi
\end{verbatim}

Finalmente, debe revisar el archivo /opt/tinyos-2.1.0/tinyos.sh tiene las siguientes líneas.

\lstinputlisting[language=bash]{tinyos_ubuntu.sh}

\subsection{Lenguaje de programación nesC}

NesC es un lenguaje de programación diseñado para sistemas embebidos que funcionan en red. NesC extiende el lenguaje C para incorporar características adicionales de concurrencia y de orientación a componentes, además modifica la estructura de un programa C para adaptarse a las restricciones de memoria de un sistema embebido. La implementación de este lenguaje realiza optimizaciones significativas para sistemas embebidos y ayuda a prevenir problemas de bajo nivel como las ''condiciones de carrera''.

\subsubsection{Componentes e interfaces}
Los principales conceptos que diferencian un programa nesC de uno C son \textbf{los componentes, las interfaces y el "wiring" (cableado)}. Estos conceptos impactan en la organización y espacio de nombres de un programa nesC. En C, los programas se separan en archivos (.c) que son enlazados mediante un espacio de nombres global, generalmente, con la ayuda de un archivo de cabecera (.h). NesC es un "lenguaje orientado a componentes", es decir, un programa nesC es una colección de componentes enlazados (cableados) vía interfaces. 

Un \textbf{componente} es un grupo de elementos relacionados funcionalmente en una unidad, similar a una clase en un lenguaje orientado a objetos. Un componente está organizado en 2 secciones: la primera sección se utiliza para especificar las interfaces que provee y utiliza, y la segunda se utiliza para implementar las funcionalidades del componente.\\

Una \textbf{interfaz} es un medio bidireccional que especifica las funciones que deben ser implementadas por el componente proveedor de la interfaz (comandos) y las funciones que deben ser implementadas por el componente usuario de la interfaz (eventos). Así, un componente que llama los comandos definidos en una interfaz, debe implementar los eventos que define esa interfaz. Un único componente puede implementar varias interfaces y múltiples instancias de la misma interfaz.\\

En nesC, existen dos tipos de componentes: los módulos y las configuraciones. Un módulo provee la implementación de una o más interfaces, mientras que las configuraciones son utilizadas para agrupar varios componentes, conectando las interfaces usadas por un componente a las interfaces provistas por otros componentes. Una aplicación nesC está finalmente descrita por una configuración de alto nivel que enlaza todos los componentes que la conforman.

En la práctica, cuando se escribe un programa nesC, se comienza con la definición de una configuración en un archivo .nc. Por convención, los archivos que describen una configuración son nombrados de la forma ''nombreC'' (por ejemplo AppC.nc). El archivo que implementa la configuración tiene la siguiente estructura.

\begin{figure}[H]
	\centering
 	\includegraphics[scale=0.5]{imgs/Ejemplo0_ModeloComponente_AppC_nc.eps} 
 	\caption{Ejemplo de configuración}
\end{figure}

La palabra reservada ''configuration'' es la que permite realizar la declaración de la configuración, en el ejemplo anterior, ''configuration AppC'' declara la configuración AppC. La palabra reservada ''implementation'' da inicio a la sección de implementación de la configuración. La sección de implementación de una configuración consiste en especificar como están conectados cada uno de los componentes que conforman la configuración. Primero se utiliza la palabra reservada ''components'' para indicar cuales son los componentes que se utilizan para construir la configuración, en el caso del ejemplo, se utilizan los componentes ficticios ''Componente1P'' y ''Componente2P''. Luego de declarar los componentes, se especifican las relaciones entre ellos utilizando los operadores $\rightarrow$ y $\leftarrow$. Estos operadores realizan un cableado directo de la interfaz que utiliza un componente, con la interfaz que provee el componente. La flecha siempre va desde el componente usuario al componente proveedor. En la figura A.7, el Componente1P provee la Interfaz1, mientras que el Componente2P utiliza la Interfaz1. El Componente2P provee la Interfaz2 y utiliza la Interfaz1\\

Las interfaces también deben ser especificadas en un archivo .nc (Por ejemplo, ''Interfaz1.nc). Un ejemplo de la estructura de un archivo que especifica una interfaz se muestra en la figura A.8.

\begin{figure}[H]
	\centering
 	\includegraphics[scale=0.5]{imgs/Ejemplo0ModeloComponenteInterfaz1nc.eps} 
 	\caption{Estructura de un archivo de ejemplo que especifica una interfaz}
\end{figure}

Para especificar una interfaz, se debe utilizar la palabra reservada ''interface'' seguida del nombre de la interfaz. El contexto que limitan las llaves se utiliza para declarar los comandos y eventos que especificará ésta interfaz.\\

El archivo que implementa el módulo Componente1P debe poseer una estructura como la que se muestra en la figura A.9.

\begin{figure}[H]
	\centering
 	\includegraphics[scale=0.5]{imgs/Ejemplo0ModeloComponenteComponente1Pnc.eps} 
 	\caption{Estructura de un archivo de ejemplo que implementa un módulo}
\end{figure}

La palabra reservada ''module'' es la que se utiliza para declarar un componente. En la primera sección, se utiliza la palabra reservada ''uses'' para declarar cuáles son las interfaces que usa el componente. También, se utiliza la palabra reservada ''provides'' para declarar cuáles son las interfaces que provee el componente. La segunda sección consiste en la implementación del módulo, donde se construyen las funcionalidades del componente, utilizando e implementando interfaces.

\subsection{Visualización de componentes}
En este tutorial se utilizará el lenguaje visual propuesto en \cite{TinyOSProgramming} para representar las relaciones entre componentes. En la figura A.10, se muestran las representaciones para los dos tipos de componentes que se utilizan en el lenguaje nesC: los módulos, son representados por cajas de bordes redondeados, dibujadas con una línea continua; las configuraciones, son representadas por cajas de bordes redondeados, dibujadas con una línea punteada. Las interfaces que provee un componente son representadas por un triángulo sólido de color negro. Si el componente provee la interfaz, el triángulo se dibuja dentro de la caja, con la base adyacente a la línea que dibuja la caja del componente. Si el componente utiliza la interfáz,  el triángulo se dibuja fuera de la caja, con la base adyacente a la línea que dibuja la caja del componente. 

\begin{figure}[h!]
	\centering
 	\includegraphics[scale=0.5]{imgs/VisualizacionDeComponentes.eps} 
 	\caption{Representación visual de componentes}
\end{figure}

En la figura A.11 se muestra la relación del ejemplo descrito en la sección anterior utilizando el lenguaje visual. 

\begin{figure}[H]
	\centering
 	\includegraphics[scale=0.5]{imgs/VisualizacionComponentesEjemplo.eps} 
 	\caption{Representación visual de una configuración}
\end{figure}

Es usual que el nombre de la configuración se coloque en el extremo superior izquierdo de la caja, para dejar espacio a las relaciones de los componentes y configuraciones que la componen. En la figura A.11 se muestra la configuración AppC, que está construida utilizando los componentes ''Componente1'' y ''Componente2''.  En este diagrama, se observa como las interfaces relacionan el Componente1 con el Componente2: El Componente1 provee la Interfaz1 y usa la Interfaz2; mientras que el Componente2 provee la Interfaz2 y usa la Interfaz1. 

\section{Introducción al Modelo de Ejecución de TinyOS}
%tares, comandos, events, async, sync, signal, splitphase.
%Overview
El modelo de ejecución de TinyOS tiene las siguientes características:
\begin{itemize}
\item Considera que todas las operaciones de entrada y salida se realizan en dos pasos (split-phase).
\item No provee hebras, en cambio, el la ejecución del programa es conducida por events (eventos) y tasks  (tareas).
\end{itemize}

En las secciones siguientes se describirán las características de las tareas, comandos y eventos.

\subsection{Tasks}
Una tarea es un mecanismo de computación diferida que permite realizar largos ciclos de ejecución y que será ejecutada hasta el final, sin importar si es interrumpida. Cuando se invoca una tarea, esta no es ejecutada de inmediato, sino que es encolada. Las tareas encoladas son ejecutadas por el Scheduler en el orden en que se encolaron y de a una a la vez. En la figura A.12 se muestra la estructura que tiene la implementación de una tarea.

\begin{figure}[H]
	\centering
 	\includegraphics[scale=0.5]{imgs/EstructuraTask1.eps} 
 	\caption{Estructura del código de una tarea.}
\end{figure}

La palabra reservada ''task'' se utiliza para declarar una tarea. Una tarea es similar a una función, puede acceder a variables, invocar comandos o funciones C, y señalizar eventos. Para invocar una tarea, se utiliza la palabra reservada ''post''. El operador post genera un valor de retorno de tipo error\textunderscore t, si la tarea no está pendiente retorna SUCCESS, si lo está retorna FAIL.\\

Las tareas, al ejecutarse una a una, no se pueden interrumpir ni adelantar entre ellas. Este esquema sincrónico de ejecución permite prescindir de herramientas como candados, semáforos o monitores.  

\subsection{Commands}
Un comando es una pieza de código que es implementada por el proveedor de una interfaz. Los comandos permiten al componente usuario de una interfaz utilizar los servicios que otro componente provee. Por ejemplo, la interfaz Leds, provee comandos para encender leds, apagar leds, alternar el estado de los leds y utilizar los leds para representar un número binario.

\begin{figure}[H]
	\centering
 	\includegraphics[scale=0.5]{imgs/interfaceLeds.eps} 
 	\caption{Especificación de la interfaz Leds}
\end{figure}

En la figura A.13 se muestra la especificación de la interfaz Leds, que se encuentra en el archivo TOS\textunderscore DIR/tos/interfaces/Leds.nc.  El comando led0On() permite encender el led0, mientras que el comando led0Off permite apagarlo. El comando get() permite obtener el estado de los leds mediante una representación binaria de sus estados, mientras que el comando set(uint8\textunderscore t val) permite establecer el estado de los leds mediante una representación binaria.\\

Se puede observar que cada declaración de comando comienza con la palabra reservada ''async''. Los comandos se ejecutan por defecto de manera sincrónica, sin embargo, nesC permite que se ejecuten comandos de manera asincrónica declarándolos como async. No es recomendable utilizar comandos o eventos asincrónicos, a menos que los tiempos sean importantes para la aplicación.\\

Como se mencionó al comienzo de la sección A.2, las operaciones de entrada y salida se realizan en dos fases, por lo que es usual que un comando involucre un callback (respuesta por parte del componente), en forma de un evento.

\subsection{Events}
Los eventos son piezas de código que deben ser implementada por el usuario de una interfaz, para manejar las respuestas a los servicios que provee un componente, que han sido iniciados por la invocación de un comando.

Un ejemplo de operación en split-phase es la lectura de un sensor utilizando la interfaz ''Read''.

\begin{figure}[H]
	\centering
 	\includegraphics[scale=0.5]{imgs/InterfaceRead.eps} 
 	\caption{Interfaz Read}
\end{figure}

Como se muestra en la figura A.14, la interfaz Read especifica el comando read() y el evento readDone(error\textunderscore t result, val\textunderscore t val), donde la llamada a read() inicia una medición y cuando la medición está lista, el componente proveedor señaliza el evento readDone, entregándole el resultado de la operación de medición (result)  y el valor obtenido de la medición (val).\\ 

Los eventos se invocan utilizando la palabra reservada ''signal''. Los eventos también son sincrónicos  por defecto , sin embargo, pueden declararse eventos asincrónicos precediendo la declaración por la palabra reservada async. 

\section{Introducción a la API de TinyOS}
Si bien TinyOS hace un esfuerzo por abstraer las interfaces y componentes del hardware, utilizando varias capaz de abstracción, gran parte de la API continúa siendo específica para cada plataforma de hardware. En esta sección, se describirán los componentes e interfaces utilizados en este trabajo para el desarrollo sobre la plataforma TmoteSky.

%UserButton y Leds
\subsection{Leds y Botón de Usuario}
Los componentes más simples del sistema son los Leds y el Botón de Usuario. Además, estos componentes son muy útiles a la hora de depurar una aplicación, obteniendo cierta realimentación con los Leds y gatillar eventos con el botón.\\

El componente UserButtonC es el que implementa el botón de usuario para el TmoteSky. Este componente, provee las interfaces Get\textless button\_state\_t\textgreater y Notify\textless button\_state\_t\textgreater.  Get permite obtener el estado del botón, mientras que Notify permite habilitar un evento que se gatilla con un cambio de estado del botón.\\

El componente LedsC es el componente que implementa los leds para el TmoteSky. Provee la interfaz Leds, que especifica una serie de comandos que permiten gestionar el estado de los leds individualmente o en conjunto.\\

Para demostrar como se pueden utilizar estos componentes, se implementa una aplicación con las siguientes funcionalidades:

\begin{itemize}
\item La aplicación debe utilizar el componente UserButtonC para monitorear los cambios de estado del botón.
\item Cada vez que se presiona el botón, la aplicación incrementa un contador y muestra una representación binaria del contador, utilizando los Leds.
\end{itemize}

Para esta aplicación se implementarán los archivos: AppC, para el código de la configuración; AppP, para el código del módulo; Makefile para el archivo de construcción de la aplicación.

\begin{figure}[H]
	\centering
 	\includegraphics[scale=0.5]{imgs/LedsBotonAppC.eps} 
 	\caption{Código de la configuración de la aplicación}
\end{figure}

En al figura A.15 se muestra el código de la configuración AppC.  En la sección de implementación, primero se declaran los componentes que se utilizarán: UserButtonC, AppP con el nombre App y LedsC. Luego se realiza el wiring de las interfaces, el componente App utilizará la implementación de la interfaz Notify que provee el componente UserButtonC, además, el componente App utilizará la implementación de la interfaz Leds que provee el componente LedsC.\\

\begin{figure}[H]
	\centering
 	\includegraphics[scale=0.5]{imgs/LedsBotonAppP.eps} 
 	\caption{Código del módulo de la aplicación}
\end{figure}

El código que se describe en la figura A.16 corresponde al del módulo de la aplicación, escrito en el archivo AppP.nc. El módulo utiliza tres interfaces: Boot, que es cableada automáticamente por el compilador y que se utiliza para implementar el código que se ejecuta cuando inicia el sistema; la interfaz Leds, que se utilizará para manejar el estado de los leds; y la interfaz Notify, que es la interfaz que provee el componente UserButtonC para monitorear los cambios de estado del botón.\\

La sección de implementación comienza con la declaración de la variable contado, que se utiliza para contar cuantas veces el usuario a presionado el botón. A continuación se implementa el evento Boot.booted(), este evento se ejecuta luego de iniciar el sistema, y permite inicializar variables y componentes. En esta aplicación, se utiliza el evento para inicializar el contador en cero y habilitar la señalización del evento Notify.notify(button\_ state\_ t b\_ state) utilizando el comando Notify.enable(). El evento Notify.notify(button\_ state\_ t b\_ state) se gatilla cada vez que hay un cambio de estado del botón, y entrega a través de la variable b\_ state, el nuevo estado. Al utilizar la interface Notify, es obligación implementar el evento Notify.notify(...), de otra forma el compilador generará como salida un error. Siempre que se utiliza una interfaz, se deben implementar todos los eventos que especifica dicha interfaz. En la implementación dele evento Notify.notify(...), se revisa que el estado entregado a través de la variable b \_ state sea igual a BUTTON\_ PRESSED, si es así se incrementa el contador y se llama al comando Leds.set(uint8\_ t val)  para mostrar a través de los led, los tres bits menos significativos del contador.\\

Los componentes, LedsC y UserButton, permiten al programador tener cierto grado de retroalimentación y de interacción en las tareas de desarrollo. En la próxima sección se describirán ejemplos para comunicar el TMoteSky con un computador, lo que permite al desarrollador disponer de herramientas aún más útiles para las tareas de desarrollo.
%Comunicación por serial
\subsection{Comunicación TMoteSky-PC}

\subsubsection{Librería printf}
Una práctica que se utiliza comúnmente para depurar programas C, es utilizar la función ''printf''. Esta función permite imprimir mensajes con valores de variables o en alguna posición determinada del código para verificar rápidamente si el programa está funcionando como el desarrollador lo espera. Para TinyOS se ha implementado una librearía que provee la función printf y que permite imprimir mensajes al puerto serial de TMoteSky. Para visualizar los mensajes es necesario ejecutar una aplicación Java en el PC, que se distribuye con el sistema, y que permite recibir los mensajes por un puerto USB e imprimirlos la consola del sistema.\\

Para utilizar la librería "printf.h" se deben realizar las siguientes modificaciones a un programa nesC.

\begin{enumerate}
\item Agregar al MAKEFILE el directorio ''TOSDIR/lib/printf'' para proveer el enlace a la librería. Opcionalmente, es posible establecer el largo del buffer utilizado para la función, por defecto tiene un tamaño de 250 bytes.

\begin{figure}[H]
	\centering
 	\includegraphics[scale=0.5]{imgs/PrintfMakefile.eps} 
 	\caption{Código del módulo de la aplicación}
\end{figure}

\item Incluir la librería ''printf.h'' en el archivo que implemente el componente que desea utilizarla. Para incluirla, basta utilizar la instrucción del pre procesador de C \# include ''printf.h''.
\end{enumerate}

Con estas modificaciones realizadas se puede utilizar la función printf(''mensaje'') para enviar mensajes por el puerto serial hacia un PC. En el lado del PC, se debe ejecutar el cliente Java para el printf que se distribuye con las herramientas Java de TinyOS. La ejecución del comando difiere dependiendo del sistema en que se está trabajando, ya que sistemas Linux y Windows los puertos seriales y USB son representados con distintos nombres. Por ejemplo en el sistema Linux Ubuntu la instrucción tiene la forma:
\begin{verbatim}
 java net.tinyos.tools.PrintfClient -comm serial@/dev/ttyUSB0:telosb
\end{verbatim}

Mientras que en un sistema windows tiene la forma:
\begin{verbatim}
 java net.tinyos.tools.PrintfClient -comm serial@COM1:telosb
\end{verbatim}

Para obtener el puerto al cuál tiene conectado el dispositivo, puede utilizar el comando ''motelist'', que se instala junto con el sistema. El parámetro ''telosb'' indica el bit rate (tasa de transmisión de bits) que se utilizará para la comunicación.

\subsubsection{Comunicación serial con SerialActiveMessageC}
TinyOS provee para la comunicación por radio el ''Radio Stack'' que provee un tipo de mensajes de largo fijo para comunicaciones single-hop denominado ActiveMessage. Para la comunicación serial TinyOS provee la misma abstracción de mensajes, lo que facilita el reenvío de los mensajes que el TmoteSky recibe por radio hacia el serial, y en sentido inverso.\\

Por el lado del PC, TinyOS prove la clase Java MoteIF, esta clase se encarga de escuchar el puerto serial o un socket, e interpretar los ActiveMessage que recibe por esa fuente y gatillar manejadores de eventos para uno o varios ActiveMessage IDs. 

%Comunicación por radio
\subsection{Comunicación por radio}
	%ActiveMessages
TinyOS provee un set de componentes e interfaces que habilitan una pila de comunicación single-hop denominada ''Radio Stack''. En la cabeza de esta pila, se encuentra la implementación ActiveMessage, que es la capa que se encarga de multiplexar el acceso de los distintos servicios de TinyOS a la radio. Esta capa utiliza el identificador ''AM Type'' para discriminar entre los servicios, similar a un puerto UDP. \\

TinyOS provee la abstracción message\_t, que corresponde a una estructura nesC que almacena punteros para una cabecera, una sección de datos, y una sección posterior a los datos para métadatos. Para acceder, modificar y enviar o recibir por el aire un message\_t, se deben utilizar las interfaces que provee el sistema, en particular, las interfaces de la capa ActiveMessage.

\begin{center}
\begin{figure}[H]
\centering
\includegraphics[scale=0.65]{imgs/EjemploAnexoRadioWiring.eps} 
\caption{Código con el cableado de los componentes necesarios para utilizar la comunicación por radio.}
\end{figure}
\end{center}

En la figura A.18 se muestran los componentes y el cableado de ellos a la aplicación que se debe realizar en el archivo de configuración para utilizar la comunicación por radio. El primer componente es ''ActiveMessageC'' que provee un control split phase para controlar la radio. En este ejemplo se cablea la interfaz split control con el nombre local de RadioControl a un módulo App, donde App es el nombre con el que se denomina al módulo principal en este ejemplo (puede ser el nombre que el programador elija). El siguiente componente utilizado es AMSender, al cuál se le asigna el nombre local EjemploSender. Este componente es instanciado con el comando new (una configuración nesC puede utilizar varias instancias de este componente con distintos nombres) y se le entrega como argumento AM\_EJEMPLO\_MSG que corresponde al AM\_TYPE que utilizará esta instancia de AMSender. AM\_TYPE es un identificador que puede asumir un valor entre 0x0000 y 0xFFFF, y su uso es similar a un puerto UDP: Permite que una aplicación utilice distintos servicios para mensajes asociados a distintos AM\_TYPE. El componente AMSender es el que provee las interfaces para construir un paquete y enviarlo por el aire utilizando la pila de comunicación ActiveMessage. El último componente utilizado es el AMReceiver, que se denomina localmente como EjemploReceiver, y es el que provee la interfaz Receive para recibir mensajes. Al igual que el componente AMSender, este componente es instanciado utilizando la palabra reservada new y con un AM\_TYPE como argumento.

\subsubsection{Definir la estructura de un mensaje} 
Para que los motes intercambien información por la radio, el primer paso es definir la estructura del mensaje que intercambiarán. No es obligación definir una estructura para los mensajes, pero tiene ventajas sobre la utilización de un buffer, ya que cuando el mensaje tiene campos con distintos tipos de datos (uint8, uint16, uint32, etc)  el código necesario para acceder a los miembros de la estructura es más corto que usar índices para acceder al buffer.  

\begin{center}
\begin{figure}[H]
\centering
\includegraphics[scale=0.65]{imgs/EstructuraMensajeEjemplo.eps} 
\caption{Ejemplo de estructura de mensaje.}
\end{figure}
\end{center}
 
Además, este tipo de estructura de mensajes permitirá utilizar la herramienta MIG. Esta herramienta, permite generar una clase Java que representa el mensaje y provee los métodos necesarios para acceder y modificar todos los campos del mensaje, en una aplicación de escritorio. 

\subsubsection{Enviar mensajes por el aire}
Para enviar un mensaje por el aire, utilizando la radio, usualmente se implementa una función que realiza las operaciones necesarias para entregar un mensaje a la pila de comunicación ActiveMessage, que se encargará de gestionar los recursos necesarios para que ese mensaje salga por el aire. A continuación, se muestra el código sugerido para enviar un mensaje del tipo Ejemplo\_Msg\_t.

\begin{center}
\begin{figure}[H]
\centering
\includegraphics[scale=0.65]{imgs/SendEjemploMsg.eps} 
\caption{Ejemplo de función que envía un mensaje del tipo Ejemplo\_Msg\_t.}
\end{figure}
\end{center}

En la figura A.20 se muestra el código necesario para enviar un mensaje por el aire. La función ''sendEjemploMsg'' se define del tipo ''error\_t'', es decir, la función sólo tiene dos posibles valores de retorno: SUCCESS o FAIL. Como argumentos recibe: destaddr, que es un entero sin signo de 16 bits que representa el TOS\_NODE\_ID del mote destinatario del mensaje; c1, campo de ejemplo 1; c2, campo de ejemplo 2; c3, campo de ejemplo 3. La función comienza con la declaración de la variable ejemploPkt, del tipo message\_t, que representa un mensaje para la pila ActiveMessage. A continuación, se revisa el flag radiobusy (se supone que esta variable existe y se utiliza para indicar si se está ocupando la radio o no): si es falso, indica que la radio está desocupada, entonces la función continúa con las operaciones de envío; si es verdadero, indica que la radio está ocupada,  entonces la función simplemente retorna FAIL. Luego, se define un puntero del tipo Ejemplo\_Msg\_t (siguiendo lo definido en la sección anterior), ejemploMsg, al cuál se le asigna la dirección del ''payload'' (carga) del mensaje message\_t declarado con anterioridad. Así, los datos que se guardan en ejemploMsg quedarán guardados en la carga del mensaje message\_t. A continuación, se asignan los valores a la estructura y posteriormente (dentro de la condición del segundo if) se realiza el envío del mensaje, llamando al comando send de la interfaz EjemploAMSend. Si el llamado retorna un valor distinto de SUCCESS, entonces la función sendEjemploMsg retorna FAIL; si el llamado retorna SUCCESS, entonces la función le asigna TRUE a radiobusy para indicar que la radio ahora estará ocupada enviando el mensaje. Como se explicó anteriormente, cada operación de entrada/salida es manejada por TinyOS en dos partes, y en este caso, solicitar a la pila de comunicación que envíe el mensaje corresponde a la primera parte. La segunda, corresponde a manejar el evento sendDone.

\begin{center}
\begin{figure}[H]
\centering
\includegraphics[scale=0.65]{imgs/EjemploSendDone.eps} 
\caption{Ejemplo de manejo del evento sendDone.}
\end{figure}
\end{center}

El código que se muestra en la figura A.21 es el que se ejecutará una vez que el mensaje ha sido enviado, en este ejemplo simplemente se vuelve a asignar FALSE al flag radiobusy.

\subsubsection{Recibir mensajes por el aire}
Para recibir un mensaje por el aire se utiliza la interfaz AMReceive. Para utilizar este componente, además de cablear el componente como se mostró al comienzo de la sección A.3.3, se debe implementar el manejador para el evento AMReceive.receive.

\begin{center}
\begin{figure}[H]
\centering
\includegraphics[scale=0.65]{imgs/EjemploReceive.eps} 
\caption{Ejemplo de manejo del evento Receive.}
\end{figure}
\end{center}

Como se muestra en la figura A.22, el evento entregará a través de los argumentos un puntero al mensaje ''msg'', un puntero a la carga del mensaje ''payload'' y el largo ''len'' del mensaje. Como se muestra en la figura, asignando ''payload'' a un puntero local del tipo de mensaje (Ejemplo\_Msg\_t) se tendrá acceso a todos los miembros de la estructura de este.

%ACKs
\subsubsection{Utilización de ACKs}
Para utilizar los ACKs se deben realizar los siguientes pasos: Primero cablear la interfaz PacketAcknowledgements del componente ActiveMessageC al módulo que utilizará los ACKs; luego, modificar la función send para que se indique en la cabecera de los mensajes que se requiere el ACK; y finalmente, revisar en el evento sendDone si se ha recibido o no el ACK.

El primer paso corresponde a cablear la interfaz PacketAcknowledgements que provee el componente ActiveMessageC al módulo que la utilizará, esto se realiza modificando el archivo de configuración ingresando una línea como se muestra a continuación.

\begin{verbatim}
MiModulo.PacketAcknowledgements->ActiveMessageC;
\end{verbatim}

Luego, se debe modificar la función send para indicar en la cabecera del mensaje que se requiere un ACK para este.

\begin{center}
\begin{figure}[H]
\centering
\includegraphics[scale=0.65]{imgs/EjemploSendAck.eps} 
\caption{Modificación a la función send para requerir ACK.}
\end{figure}
\end{center}

En la figura A.23 se muestra la función send modificada para requerir ACK para un mensaje. La modificación consiste en agregar una línea para invocar el comando requestAck de la interfaz PacketAcknowledgements. Luego, para verificar la recepción de un ACK se debe modificar el evento sendDone como se muestra a continuación.

\begin{center}
\begin{figure}[H]
\centering
\includegraphics[scale=0.65]{imgs/EjemploSendDoneAck.eps} 
\caption{Modificación al evento sendDone revisar recepción de ACK.}
\end{figure}
\end{center}

En la figura A.24 se muestra la modificación para revisar si el ACK fue recibido. La modificación consiste en invocar el comando wasAcked() de la interfaz PacketAcknowledgements.

%RSSI
\subsubsection{Obtener el RSSI de un mensaje recibido}
Para obtener el RSSI de un mensaje recibido es necesario modificar el archivo con la configuración para agregar el componente necesario. Se debe agregar la siguiente línea a la declaración de componentes.

\begin{verbatim}
components CC2420ActiveMessageC;
\end{verbatim}

A continuación, se debe cablear la interfaz CC2420Packet al módulo.

\begin{verbatim}
MiModulo->CC2420ActiveMessageC.CC2420Packet;
\end{verbatim}

Finalmente, se debe modificar el evento receive como se muestra a continuación.

\begin{center}
\begin{figure}[H]
\centering
\includegraphics[scale=0.65]{imgs/EjemploSendDoneAck.eps} 
\caption{Modificación al evento receive para obtener el RSSI de un mensaje.}
\end{figure}
\end{center}

%SensoresTmoteSky
%\subsubsection{Utilización del sensor de Luz}

%LowPowerListening
\subsubsection{Operación de bajo consumo - Low Power Listening}
Para aplicaciones que requieran extender el tiempo de duración de la carga de las baterías, TinyOS provee la herramienta Low Power Listening. Esta herramienta, reduce el consumo de energía disminuyendo el tiempo de ''idle listening''. Para lograr esto, se implementa un ciclo de trabajo para la radio donde primero se activa la recepción por el tiempo necesario para detectar un ''carrier'' en el canal; si no hay carrier, apaga la radio por un intervalo de tiempo denominado ''Intervalo LPL''; si hay carrier, se mantiene la radio encendida para recibir el mensaje. Por otro lado, para lograr entregar un mensaje, el transmisor extiende el preámbulo del mensaje hasta que la duración temporal alcance el intervalo LPL. En el anexo C se profundiza en la herramienta Low Power Listening.\\

Para utilizar este componente, primero se deben declarar y cablear los siguientes componentes en el archivo de configuración.

\begin{verbatim}
components CC2420ActiveMessageC;
MiModulo.LowPowerListening -> CC2420ActiveMessageC;
\end{verbatim}

Luego, para iniciar el ciclo de trabajo en el receptor, se debe invocar el siguiente comando.
\begin{center}
\begin{figure}[H]
\centering
\includegraphics[scale=0.65]{imgs/InvocacionLowPowerListening.eps} 
\caption{Invocación para iniciar el ciclo de trabajo de Low Power Listening.}
\end{figure}
\end{center}

En el transmisor, se debe agregar la siguiente modificación a la función de envío del mensaje.

\begin{center}
\begin{figure}[H]
\centering
\includegraphics[scale=0.65]{imgs/SendLPL.eps} 
\caption{Modificación al código de la función send para utilizar LPL.}
\end{figure}
\end{center}

La modificación que se muestra en la figura A.27 consiste en agregar la línea que invoca al comando ''LowPowerListening.setRxSleepInterval(\&ejemploPkt). Este comando marca el paquete para que la pila de comunicación extienda el preámbulo del envío para que el receptor que utiliza LPL lo pueda recibir correctamente.\\

Finalmente, es necesario modificar el Makefile de la aplicación para que el compilador utiliza la pila de comunicación que implementa LPL para construir el archivo binario. La modificación consiste en agregar el siguiente flag de compilación.
\begin{center}
\begin{verbatim}
CFLAGS += -DLOW_POWER_LISTENING
\end{verbatim}
\end{center}

Con estas modificaciones la aplicación debe habilitar correctamente las operaciones de comunicación de bajo consumo.

    %\
    %Incluir anexo B
    %	
\chapter{Implementación de WSNSniffer}
Para evaluar el comportamiento lógico de la red de sensores inalámbricos, fue necesario contar con una aplicación que permitiera visualizar los mensajes intercambiados  entre los  nodos de la red. Lamentablemente, no se ha publicado una aplicación de escritorio que permita realizar esto, sin embargo, junto con TinyOS se distribuyen las herramientas necesarias para desarrollar un sistema que cumpla con los siguientes objetivos.

\begin{itemize}
 \item Capturar los mensajes intercambiados por nodos de la red.
 \item Enviar esos mensajes por el puerto USB del TmoteSky.
 \item Presentar los mensajes en una aplicación de escritorio, que utilice una interfaz gráfica de usuario.
 \item Generar reportes en forma de tablas excel, para contar con resultados persistentes.
\end{itemize}

Utilizando las herramientas que provee TinyOS para el lenguaje de programación Java, se implementó una aplicación de escritorio que cumpliera con los objetivos descritos anteriormente. Además, se modificó la aplicación nesC ''BaseStation'' para que escuchara todos los mensajes que se hay en el canal y se instaló en un TMoteSky para que alimentara de mensajes a la aplicación Sniffer de escritorio. Ambas aplicaciones conforman el sistema que se muestra en la figura B.1
\begin{figure}[H]
	\centering
 	\includegraphics[scale=1]{imgs/ArquitecturaSistemaSniffer.eps} 
 	\caption{Arquitectura del Sistema Sniffer}
\end{figure}

En las próximas secciones se describen las etapas de este sistema implementado.

\section{Etapa de captura y envío por USB de los mensajes detectados en el canal}
Dentro de las aplicaciones nesC que incluye TinyOS se encuentra la aplicación "BaseStation". Esta aplicación funciona como un puente entre los mensajes que el mote recibe por el aire y un PC. Como el puerto serial está conectado a un microcontrolador USB, los mensajes pueden ser leídos por un PC a través de este puerto.  Por defecto, la aplicación "BaseStation" sólo escucha los mensajes que son direccionados al TOS\textunderscore NODE\textunderscore ID del nodo. Para habilitar que escuche todos los mensajes en el canal se debe modificar el archivo Makefile de la aplicación agregar el flag -DCC2420\textunderscore NO\textunderscore ADDRESS \textunderscore RECOGNITION. De esta forma, queda implementada la primera etapa del sistema Sniffer.

\section{Etapa de presentación del detalle de los mensajes}
Para implementar una aplicación que presente los mensajes recibidos de la red se opta por desarrollar una aplicación de escritorio Java. Las razones de esta elección son las siguientes:

\begin{itemize}
\item TinyOS incluye herramientas programadas en este lenguaje para interactuar con un nodo por comunicación serial.
\item Permite el desarrollo rápido de interfaces gráficas de usuario.
\item Permite la integración rápida con bases de datos y APIs para generar archivos Excel.
\end{itemize}

\begin{figure}[H]
	\centering
 	\includegraphics[scale=0.30]{imgs/SnifferScreenshot.eps} 
 	\caption{Captura de pantalla de la aplicación WSNSniffer}
\end{figure}

Se utiliza el IDE Netbeans para crear una aplicación con una interfaz gráfica que consta de una pequeña barra de menú y tres pestañas: la primera pestaña muestra la vista ''BaseStation'', que está pensada para ser usada con un mote utilizando la aplicación ''BaseStation'' y recibe los mensajes direccionados al Gateway y muestra, en función del nodo, el dato medido actualizado con el último valor recibido; en la segunda pestaña se muestra la vista ''Sniffer'', pensada para ser utilizada con un mote corriendo la aplicación Sniffer, que muestra todos los mensajes captados en el canal de comunicación; finalmente, la última pestaña muestra estadísticas sobre los mensajes recibidos. Cada pestaña utiliza una JTable para mostrar la información correspondiente.

\section{Recepción de mensajes por el puerto serie}
%intro
Como se mencionó antes, la aplicación BaseStation en el TmoteSky, reenvía los mensajes recibidos por la radio a través de la UART del MSP430, que mediante un microcontrolador USB, son enviados al puerto USB. Los mensajes reenviados por el puerto serie son compatibles con la estructura ActiveMessages, que es la misma que se utiliza para intercambiar mensajes por radio, por lo que los mensajes que llegan por el puerto serie a la aplicación  son exactamente los mismos que recibe el mote por radio.\\

% La pila ActiveMessage implementa la generación de ACKs por mensajes recibidos, pero estos mensajes no se pueden capturar, ya que los ACKs utilizados son de bajo nivel y no se proveen funciones para obtener estos mensajes en las capas de alto nivel.\\

%La clase MoteIF
Para recibir un ActiveMessage por el puerto USB, la colección de herramientas Java de TinyOS provee la interfaz a nivel de aplicación "net.tinyos.message.MoteIF". Esta clase recibe de la aplicación un  \textbf{packet source}, que es un medio de comunicación por el cuál la clase MoteIF intercambia paquetes ActiveMessage con TinyOS corriendo en el TMoteSky. Un packet source puede ser un puerto serial, un socket TCP o la herramienta SerialForwarder. Para el desarrollo de esta aplicación, se opta por utilizar como packet source el puerto serial.\\

%Manejo de eventos, registrar listeners.
La clase MoteIF, permite registrar "listeners" (aclaración) para mensajes ActiveMessage de un "MessageType" en particular. Así, utilizando un manejador de eventos, se puede gestionar la información de un mensaje recibido. Dada esta estructura de recepción de mensajes, la aplicación sólo puede gestionar los mensajes de "MessageType" conocidos, es decir, que el programador sabe cuáles son los "Message Type" que se están intercambiando en el canal. En consecuencia, si no se registra un Listener, la clase MoteIF ignorará el mensaje. El método MoteIF.RegisterListener, recibe como entrada una clase Message o una clase heredera de esta y una implementación de la interfaz MessageListener. La clase que implementa la interfaz MessageListener, debe implementar el método void messageReceived(int to, Message m), que es el método que maneja el evento de recepción del mensaje al MessageType asociado.\\

%Clases de mensaje y Mig
En esta aplicación se utilizan dos implementaciones de la interfaz MessageListener: TopologyMsgListener y MHDataMsgListener, que corresponden a las dos clases de mensajes que gestiona la aplicación: TopologyMsg y MHDataMsg. Las clases de mensajes implementadas se ocupan de estructurar los mensajes recibidos, es decir, entregan los métodos get y set necesarios para cada campo del mensaje, una sobrecarga del método toString, que permite crear un string con toda la información que describe el mensaje y los métodos para crear nuevos mensajes desde la aplicación. Estas clases son en parte creadas con la herramienta "MIG" que se distribuye con TinyOS. MIG (Message Interface Generator) es una herramienta que genera código Java, Python o C para gestionar los ActiveMessage declarados en el archivo de cabecera ".h" que contiene las estructuras  que utilizan las aplicaciones nesC de los motes.\\

Sin embargo, la herramienta MIG no puede generar la clase requerida para el tipo de mensaje MHDataMsg, debido a que el largo de este tipo de mensajes no es fijo. Para solucionar este problema, se utiliza MIG para generar la cabecera de los mensajes y el manejo de el campo "data", de largo variable, se implementa dentro de la misma clase pero de forma manual. El comando que se utiliza para generar una clase java a partir de la descripción de un mensaje en un .h es:\\

\begin{verbatim}
mig -target=telosb -java-classname=TopologyMsg java Messages.h 
Topology_Msg -o TopologyMsg.java

mig -target=telosb -java-classname=DebugMsg java Messages.h 
Debug_Msg -o DebugMsg.java
\end{verbatim}

Luego de capturar los mensajes, estos se ingresan a una JTable para presentarlos visualmente. Se asignan colores a los mensajes de topología, para diferenciarlos claramente de los mensajes de datos en las pruebas. 

\section{Generación de archivos Excel}
La cantidad de mensajes que puede escuchar el sniffer en una prueba puede ir desde decenas hasta millones. Cuando una aplicación está presentando millones de mensajes, es posible la aplicación se vuelva inestable y se arriesga a perder los datos recolectados. Por lo tanto, es necesario pasar esa información a otro medio. Se estudia la posibilidad de incorporar una base de datos a la aplicación y también, la posibilidad de generar archivos Excel copiando la tabla presentada en la aplicación a un archivo permanente. Debido a que la base de datos requiere mucho más desarrollo y que la generación de algún tipo de archivo o reporte será necesario en cualquier caso, se opta por utilizar una API que genere archivos Excel.\\

Para generar un archivo Excel, se utiliza la Java Excel API, que es una librería de código abierto que permite leer, escribir, modificar y crear planillas Excel. Esta API se distribuye en el archivo jxl.jar. Para que el usuario pueda generar un archivo Excel, se agrega un botón en el menú archivo, que ejecuta la rutina para exportar la tabla a un archivo Excel.\\

En el CD adjunto a este trabajo, se encuentra el proyecto Netbeans que implementa esta aplicación. En el se encuentra el código fuente y también las versión jar. Además, se incluye también la versión del WSNSniffer para el protocolo PEGASIS. Es importante destacar que para que estos programas funcionen correctamente, es necesario tener correctamente instaladas las herramientas de TinyOS que permiten la comunicación serial con los motes.

    
    %Incluir anexo C
    %	\chapter{ Low Power Listening}
El componente Low Power Listening (LPL) es una herramienta que TinyOS provee para el desarrollo de aplicaciones de bajo consumo. La principal característica de LPL es que permite reducir el consumo energético manteniendo la comunicación asincrónica. LPL reduce el consumo de energía de un mote atacando una de las acciones que desperdicia más energía en un sistema de comunicaciones, el ''Idle Listening''. En el Anexo A, sección A.3.3 se describen las modificaciones que se deben realizar a una aplicación para incorporar el componente LPL.\\

Para reducir el consumo, el componente LPL aplica un ciclo de trabajo al hardware de comunicación: la radio se deshabilita por un tiempo denominado ''Intervalo LPL''; una vez que ha transcurrido el intervalo, se enciende sólo el tiempo necesario para detectar si es que hay algún otro mote transmitiendo, es decir, si hay un carrier en el canal; si hay, la radio se enciende y el mote recibe el mensaje; si no hay, se deshabilita la radio nuevamente. Para que la comunicación asincrónica se realice correctamente, el mote transmisor  extiende el preámbulo del mensaje por al menos la duración del intervalo LPL, para así garantizar que el mote receptor a revisado el canal al menos una vez desde que se comenzó hasta que se terminó la operación de envío.

\begin{figure}[H]
 \centering
 \includegraphics[scale=0.35]{imgs/LplNoLpl.eps}
 \caption{Corriente medida en un TMoteSky con y sin utilización del componente LPL}
\end{figure}

En la figura C.1 se muestra una comparación entre el comportamiento de la corriente utilizada por el TmoteSky cuando no utiliza LPL (en rojo) y cuando utiliza (en azul). En esta medición, se utilizó un intervalo LPL de 50[ms] y en la figura se puede observar como cambia el nivel de corriente utilizado por el mote al operar con el componente LPL. Cuando el mote suspende la radio, se observa que el nivel de corriente utilizada baja a menos de 1[mA] en promedio, y cuando habilita la radio nuevamente para revisar el canal el nivel de corriente sube hasta 17[mA] aproximadamente.\\

De acuerdo con el principio de funcionamiento del LPL, el consumo energético debería variar principalmente por: la duración del intervalo LPL; la cantidad de mensajes recibidos por un mote; y la cantidad de mensajes enviados. Para analizar esto, se realizan dos experimentos: el primer experimento consiste en medir la corriente media utilizada por un mote, en función del intervalo LPL configurado en la aplicación, sin ningún mote transmitiendo; el segundo consiste en medir la corriente media utilizada en dos motes, uno que transmite mensajes y otro que recibe mensajes, con distintos intervalos LPL. Del primer experimento se obtienen los resultados que se observan en la figura C.2.

\begin{figure}[H]
 \centering
 \includegraphics[scale=0.5]{imgs/CaracLplConsumoVsIntervalo.eps}
 \caption{Gráfico de la corriente media utilizada por el TmoteSky en función el intervalo LPL}
\end{figure}

Como se observa en la figura C.2, se realizan mediciones de la corriente media para intervalos LPL que van desde los 50[ms] hasta 1000[ms], incrementándolos en 50[ms] para cada medición. En la figura se observa que a medida que se incrementa el intervalo LPL, la corriente media disminuye: para un intervalo LPL de 50 [ms], resulta una corriente media de 3.15[mA]; para un intervalo LPL de 500[ms], la corriente media resultante fue 0.7558[mA]; para un intervalo LPL de 1000[ms], la corriente media resultante fue 0.67[mA]. La disminución no se produce de manera proporcional, sino que más bien, parece disminuir exponencialmente. Según los resultados obtenidos, utilizando un intervalo LPL de 50[ms] y dos baterías alcalinas AA, un TMoteSky podría funcionar por 31 días. Con la misma batería y un intervalo LPL de 500[ms], el dispositivo podría funcionar por 132 días. Utilizando dos baterías alcalinas D con un intervalo LPL de 500[ms], el dispositivo podría funcionar por 661 días. Con la misma batería y utilizando un intervalo LPL de 1000[ms], el dispositivo podría funcionar por 741 días. Sin embargo, esto experimento no contempla envío ni recepción de mensajes por parte del dispositivo, operaciones que aumentan considerablemente el consumos de energía.\\

%Segundo Experimento
Para medir el impacto energético de las operaciones de envío y recepción de mensajes, se implementan dos aplicaciones de prueba, una que recibe mensajes y otra que envía mensajes periódicamente, ambas utilizando el componente LPL. El segundo experimento consiste en medir la corriente consumida por un TMoteSky que envía mensajes periódicamente utilizando LPL (con el dispositivo receptor funcionando) y, posteriormente, medir el consumo en el dispositivo receptor.

\begin{figure}[H]
 \centering
\includegraphics[scale = 0.5]{imgs/CaractLplCorrienteTx.eps}
\caption{Corriente media utilizada por el transmisor de un mensaje para distintos intervalos de envío de mensajes y distintos intervalos LPL }
\end{figure}

La figura C.3 muestra un gráfico que corresponde a la corriente media utilizada por el dispositivo transmisor para tres intervalos LPL. Para cada intervalo LPL, se utilizaron además tres intervalos para la generación de mensajes: 100, 500 y 1000 milisegundos. Se observa que para todos los intervalos LPL, la corriente media medida disminuye al aumentar el intervalo de generación de paquetes. Además, se observa que la corriente media utilizada aumenta al aumentar la duración del intervalo LPL. Esto se debe al principio del funcionamiento del componente LPL, ya que este obliga al transmisor a mantener la radio encendida un tiempo igual o mayor al de duración del intervalo LPL del receptor, con el uso intensivo de corriente que esto conlleva.

\begin{figure}[H]
 \centering
\includegraphics[scale = 0.5]{imgs/CaractLplCorrienteRx.eps}
\caption{Corriente media consumida por el receptor de un mensaje para distintos intervalos de generación de mensajes y distintos intervalos LPL}
\end{figure}

En la figura C.4, se muestra un gráfico que corresponde a la corriente media utilizada por el receptor. Se observa que la corriente media, en función  del intervalo de generación de mensajes, tiene el mismo comportamiento que en el transmisor, es decir, disminuye a medida que se aumenta la duración del intervalo. Sin embargo, se observa que la corriente media en función el intervalo LPL tiene el comportamiento inverso al que se observa en el transmisor, es decir, la corriente media disminuye a medida que se incrementa la duración del intervalo LPL.

Con estos experimentos, se observa claramente que existe un compromiso en la elección de un intervalo LPL: un intervalo muy corto aumenta la corriente media utilizada para la operación de escuchar mensajes, pero mantiene las operaciones de transmitir con una corriente media baja; un intervalo muy largo, disminuye la corriente media utilizada para escuchar mensajes, pero incrementa la corriente media utilizada para las operaciones de transmisión. Por otro lado, se observa que a medida que se aumenta la duración del intervalo de generación de mensajes, la corriente media utilizada siempre disminuye, por lo que para disminuir el consumo de energía de un mote se debe utilizar el intervalo de generación de mensajes más largo que permita la aplicación. Finalmente, se observa que LPL puede ayudar a reducir la corriente utilizada por los motes a menos de 3 [mA], que permite una autonomía de los motes del orden de meses (utilizando como referencia una batería AA), y manteniendo la comunicación asincrónica. Observando los resultados obtenidos, se plantea que para el futuro se puede desarrollar un modelo matemático para estimar el consumo de un mote, con sólo conocer ciertos parámetros del funcionamiento de la aplicación y del componente LPL; además, se debe investigar como es afectado el nivel de corriente utilizado por un mote frente a tráfico ajeno en el canal (tráfico no direccionado, ni enviado por el mote).
    
\end{document}
